
TFHE 方案用于整数的同态加法与乘法，尤其适合逻辑电路这样的按位计算。与 BFV、BGV 或 CKKS 不同，TFHE 的特点是噪声自举速度快，因此能够高效处理乘法深度较大的计算。TFHE 的噪声自举技术还可以进一步用于函数加密。

在 TFHE 中，每个明文都被加密为一个 LWE 密文。因此，TFHE 的密文与密文加法、密文与明文加法，以及密文与明文乘法，均基于 \autoref{part:generic-fhe} 中介绍的 GLWE 同态加法和乘法实现；取 $n = 1$ 即可将 GLWE 化为 LWE。

本节介绍 TFHE 的特有组成部分：密钥切换、密文与密文乘法、系数提取，以及噪声自举。

$ $

\begin{tcolorbox}[
    title = \textbf{所需背景知识},    % box title
    colback = white,    % light background; tweak to taste
    colframe = black,  % frame colour
    boxrule = 0.8pt,     % line thickness
    left = 1mm, right = 1mm, top = 1mm, bottom = 1mm % inner padding
]
\begin{itemize}
  \item \autoref{sec:modulo}: \nameref{sec:modulo}
  \item \autoref{sec:group}: \nameref{sec:group}
  \item \autoref{sec:field}: \nameref{sec:field}
  \item \autoref{sec:order}: \nameref{sec:order}
  \item \autoref{sec:polynomial-ring}: \nameref{sec:polynomial-ring}
  \item \autoref{sec:decomp}: \nameref{sec:decomp}
  \item \autoref{sec:modulus-rescaling}: \nameref{sec:modulus-rescaling}
  \item \autoref{sec:lattice}: \nameref{sec:lattice}
  \item \autoref{sec:lwe}: \nameref{sec:lwe}
  \item \autoref{sec:rlwe}: \nameref{sec:rlwe}
  \item \autoref{sec:glwe}: \nameref{sec:glwe}
  \item \autoref{sec:glev}: \nameref{sec:glev}
  \item \autoref{sec:ggsw}: \nameref{sec:ggsw}
  \item \autoref{sec:glwe-add-cipher}: \nameref{sec:glwe-add-cipher}
  \item \autoref{sec:glwe-add-plain}: \nameref{sec:glwe-add-plain}
  \item \autoref{sec:glwe-mult-plain}: \nameref{sec:glwe-mult-plain}
  \item \autoref{subsec:modulus-switch-lwe}: \nameref{subsec:modulus-switch-lwe}
  \item \autoref{sec:glwe-key-switching}: \nameref{sec:glwe-key-switching}
\end{itemize}
\end{tcolorbox}

\clearpage

\subsection{加密与解密}
\label{subsec:tfhe-enc-dec}

TFHE 基于 LWE 密码系统（\autoref{sec:lwe}）进行加密和解密，它等价于取 $n = 1$ 的 GLWE 密码系统（\autoref{sec:glwe}）。不过，与前面介绍的 LWE 密码系统不同，TFHE 从二元集合 $\{0, 1\}$ 中采样秘密密钥的元素，而不是从三元集合 $\{-1, 0, 1\}$ 中采样。


\begin{tcolorbox}[title={\textbf{\tboxlabel{\ref*{subsec:tfhe-enc-dec}} TFHE 加密与解密}}]
\textbf{\underline{初始设置}：} $\Delta = \dfrac{q}{t}$, $\vec{s} \xleftarrow{\$} \mathbb{Z}_2^k$  \textcolor{red}{ $\rhd$ 其中 $t$ 整除 $q$，且 $\vec{s}$ 的每个元素都是零次多项式}

$ $

\par\noindent\rule{\textwidth}{0.4pt}

\textbf{\underline{加密输入}：} $m \in \mathbb{Z}_t$, $\vec{a} \xleftarrow{\$} \mathbb{Z}_q^k$, $e \xleftarrow{\chi_\sigma} \mathbb{Z}_q$ \textcolor{red}{ $\rhd$ $\vec{a}$ 的每个元素都是零次多项式}
\begin{enumerate}
\item 放大尺度：$m \longrightarrow \Delta \cdot m \text{ } \in \mathbb{Z}_q$

\item 计算 $b = \vec{a} \cdot \vec{s} + \Delta  m + e \text{ } \bmod q$
\item $\textsf{LWE}_{\vec{s},\sigma}(\Delta  m + e) = (\vec{a}, b) \text{ } \in \mathbb{Z}_q^{k + 1}$ 
\end{enumerate}

\par\noindent\rule{\textwidth}{0.4pt}

\textbf{\underline{解密输入}：} $\textsf{ct} = (\vec{a}, b) \text{ } \in \mathbb{Z}_q^{k+1}$
\begin{enumerate}
\item $\textsf{LWE}^{-1}_{\vec{s},\sigma}(\textsf{ct}) = b - \vec{a}\cdot \vec{s} = \Delta  m + e  \pmod q$

\item 缩小尺度：$\Bigg\lceil\dfrac{ \Delta  m + e } {\Delta}\Bigg\rfloor = m \text{ } \in \mathbb{Z}_t$ \textcolor{red}{ $\rhd$ 即执行 $q \rightarrow t$ 的模数切换}
\end{enumerate}

$ $

\textbf{{正确解密的条件}：}
\begin{itemize}
\item 同态运算过程中增长的噪声 $e$ 应满足：$e < \dfrac{\Delta}{2}$。
\end{itemize}

\end{tcolorbox}

为简化记号，本节经常将 $\textsf{LWE}_{\vec{s},\sigma}(\Delta  m + e)$ 写成 $\textsf{LWE}_{\vec{s},\sigma}(\Delta  m)$，因为 $\textsf{LWE}_{\vec{s},\sigma}(\Delta  m + e) \approx \textsf{LWE}_{\vec{s},\sigma}(\Delta  m)$（即二者解密得到同一消息）。即使我们用 $\textsf{LWE}_{\vec{s},\sigma}(\Delta  m)$ 代替 $\textsf{LWE}_{\vec{s},\sigma}(\Delta  m + e)$，也应将其理解为对 $\Delta m + e$ 的加密，即放大后的消息中仍然包含噪声。

\subsection{密文与密文的同态加法}
\label{subsec:tfhe-add-cipher}

TFHE 的密文与密文加法采用 LWE 的密文加法方案，它等价于取 $n = 1$ 的 GLWE 密文加法方案（\autoref{sec:glwe-add-cipher}）。


\begin{tcolorbox}[title={\textbf{\tboxlabel{\ref*{subsec:tfhe-add-cipher}} TFHE 密文与密文加法}}]
$\textsf{LWE}_{\vec{s}, \sigma}(\Delta m^{\langle 1 \rangle} + e^{\langle 1 \rangle} ) + \textsf{LWE}_{\vec{s}, \sigma}(\Delta m^{\langle 2 \rangle} + e^{\langle 2 \rangle}) $

$ = ( \vec{a}^{\langle 1 \rangle}, \text{ } b^{\langle 1 \rangle}) + (\vec{a}^{\langle 2 \rangle}, \text{ } b^{\langle 2 \rangle}) $

$ = ( \vec{a}^{\langle 1 \rangle} + \vec{a}^{\langle 2 \rangle}, \text{ } b^{\langle 1 \rangle} + b^{\langle 2 \rangle} ) $

$= \textsf{LWE}_{\vec{s}, \sigma}(\Delta(m^{\langle 1 \rangle} + m^{\langle 2 \rangle}) + e^{\langle 1 + 2 \rangle} )$
\label{Here}
\end{tcolorbox}


\subsection{密文与明文的同态加法}
\label{subsec:tfhe-add-plain}

TFHE 的密文与明文加法（其中 $\lambda$ 是待加的常数）采用 LWE 的密文与明文加法方案，它等价于取 $n = 1$ 的 GLWE 密文与明文加法方案（\autoref{sec:glwe-add-plain}）。

\begin{tcolorbox}[title={\textbf{\tboxlabel{\ref*{subsec:tfhe-add-plain}} TFHE 密文与明文加法}}]
$\textsf{LWE}_{\vec{s}, \sigma}(\Delta m + e) + \Delta \lambda $

$=  (\vec{a}, \text{ } b) + \Delta \lambda$

$=  (\vec{a}, \text{ } b + \Delta\lambda)$

$= \textsf{LWE}_{\vec{s}, \sigma}(\Delta (m + \lambda) + e )$
\end{tcolorbox}


\subsection{密文与明文的同态乘法}
\label{subsec:tfhe-mult-plain}

TFHE 的密文与明文乘法采用 LWE 的密文与明文乘法方案，它等价于取 $n = 1$ 的 GLWE 密文与明文乘法方案（\autoref{sec:glwe-mult-plain}）。

\begin{tcolorbox}[title={\textbf{\tboxlabel{\ref*{subsec:tfhe-mult-plain}} TFHE 密文与明文乘法}}]
$\textsf{LWE}_{\vec{s}, \sigma}(\Delta m + e) \cdot \lambda$

$= (\vec{a}, \text{ } b) \cdot \lambda$

$= (\lambda\cdot \vec{a}, \text{ } \lambda \cdot b)$

$= \textsf{LWE}_{\vec{s}, \sigma}(\Delta (m \cdot \lambda) + \lambda \cdot e )$
\end{tcolorbox}


\subsection{同态密钥切换}
\label{subsec:tfhe-key-switching}

\textbf{- 参考资料：}
\href{https://www.zama.ai/post/tfhe-deep-dive-part-3}{深入理解 TFHE：第三部分——密钥切换与分层乘法}~\cite{tfhe-3}

TFHE 的密钥切换方案将 LWE 密文对应的秘密密钥从 $\vec{s}$ 变为 $\vec{s}_{'}$，两个密钥向量的维数可以相同，也可以不同。该方案本质上就是 LWE 密钥切换。具体而言，它等价于取 $n = 1$ 的另一种 GLWE 表示（\autoref{subsec:glwe-alternative}）所采用的密钥切换方案（\autoref{sec:glwe-key-switching}），如下所示：

\begin{tcolorbox}[title={\textbf{\tboxlabel{\ref*{subsec:tfhe-key-switching}} TFHE 密钥切换}}]

给定 $\textsf{LWE}_{\vec{s},\sigma}(\Delta m + e) = (\vec{a}, b)$，

$\textsf{LWE}_{\vec{s}_{'},\sigma}(\Delta m + e') = (0, b) - \bm{\langle} \textsf{Decomp}^{\beta, l}(\vec{a}), \text{ } \textsf{Lev}_{\vec{s}_{'}, \sigma}^{\beta, l}(\vec{s}) \bm{\rangle}$
\end{tcolorbox}


\subsection{密文与密文的同态乘法}
\label{subsec:tfhe-mult-cipher}

\textbf{- 参考资料：}
\href{https://www.zama.ai/post/tfhe-deep-dive-part-3}{深入理解 TFHE：第三部分——密钥切换与分层乘法}~\cite{tfhe-3}

$ $

TFHE 支持以下形式的两个密文相乘：$\textsf{LWE}_{\vec{s}, \sigma}(\Delta m_1) \cdot \textsf{GSW}_{\vec{s}, \sigma}^{\beta, l}(m_2)$。

$ $

\noindent 第一项 $\textsf{LWE}_{\vec{s}, \sigma}(\Delta m_1 + e_1)$ 可由以下任一种方式得到：
\begin{itemize}
\item 对明文 $m_1$ 进行一次新的 LWE 加密（\autoref{subsec:glwe-enc}）。
\item 对两个 LWE 密文执行同态加法所得的结果（\autoref{sec:glwe-add-cipher}）。
\item 对一个 LWE 密文和一个明文执行同态乘法所得的结果（\autoref{sec:glwe-mult-plain}）。
\end{itemize}

$ $

\noindent 第二项 $\textsf{GSW}_{\vec{s}, \sigma}^{\beta, l}(m_2)$ 可由以下任一种方式得到：
\begin{itemize}
\item 对明文 $m_2$ 进行一次新的 GSW 加密（\autoref{subsec:ggsw-enc}）。
\item 通过\textit{电路自举}将 $\textsf{LWE}_{\vec{s}, \sigma}(\Delta m_2 + e_2)$ 转换为 $\textsf{GSW}_{\vec{s}, \sigma}^{\beta, l}(m_2)$（将在未来版本中介绍）。
\end{itemize}

$ $

\noindent 回顾以下关系：

$\textsf{LWE}_{\vec{s}, \sigma}(\Delta m_1 + e_1) = (\vec{a}, b) \in \mathbb{Z}_{q}^{k + 1}$，其中 $b = \vec{a} \cdot \vec{s} + \Delta m_1 + e_1$

$ $

$\textsf{GSW}_{\vec{s}, \sigma}^{\beta, l}(m_2) = \Bigr \{ \{ \textsf{Lev}_{\vec{s}, \sigma}^{\beta, l} (-s_i \cdot m_2)  \}_{i=0}^{k-1}, \textsf{Lev}_{\vec{s}, \sigma}^{\beta, l}(m_2) \Bigl \} \in \mathbb{Z}_{q }^{(k+1) \cdot l \cdot (k'+1)}$ \textcolor{red}{ $\rhd$ 由 \autoref{subsec:ggsw-enc} 可得}

$ $

\noindent 使用以下记号：

$\textsf{GSW}_{\vec{s}, \sigma}^{\beta, l}(m_2) = {\bar{\textsf{ct}}} = (\bar{\textsf{ct}}_0,. \bar{\textsf{ct}}_1, \gap{$\cdots$} \bar{\textsf{ct}}_k)$ 

$\bar{\textsf{ct}}_i = \textsf{Lev}_{\vec{s}, \sigma}^{\beta, l}(-s_i \cdot m_2)$，其中 $0 \leq i \leq (k-1)$

$\bar{\textsf{ct}}_k = \textsf{Lev}_{\vec{s}, \sigma}^{\beta, l}(m_2)$

$\textsf{ct} = \textsf{LWE}_{\vec{s}, \sigma}(\Delta m_1 + e_1) = (\vec{a}, b) = (a_0, a_1, \gap{$\cdots$}, a_{k-1}, b) = (\textsf{ct}_0, \textsf{ct}_1, \cdots, \textsf{ct}_k)$

$ $

\noindent 定义以下 TFHE 密文乘法运算：

$\textsf{ct} \cdot {\bar{\textsf{ct}}} = \sum\limits_{i=0}^{k}\langle \textsf{Decomp}^{\beta, l}(\textsf{ct}_i), \bar{\textsf{ct}}_i \rangle$

$ $

\noindent 于是下式成立：

\begin{tcolorbox}[title={\textbf{\tboxlabel{\ref*{subsec:tfhe-mult-cipher}} TFHE 密文与密文乘法}}]
$\textsf{ct} = \textsf{LWE}_{\vec{s}, \sigma}(\Delta m_1 + e_1) = (a_0, a_1, \cdots, a_{k-1}, b)$

$\bar{\textsf{ct}} = \textsf{GSW}_{\vec{s}, \sigma}^{\beta, l}(m_2) = \bm( \textsf{Lev}_{\vec{s}, \sigma}^{\beta, l}(-s_0\cdot m_2), \textsf{Lev}_{\vec{s}, \sigma}^{\beta, l}(-s_1\cdot m_2), \cdots, \textsf{Lev}_{\vec{s}, \sigma}^{\beta, l}(-s_{k-1}\cdot m_2), \textsf{Lev}_{\vec{s}, \sigma}^{\beta, l}(m_2)  \bm)$

$\textsf{LWE}_{\vec{s}, \sigma}(\Delta m_1 + e_1) \cdot \textsf{GSW}_{\vec{s}, \sigma}^{\beta, l}(m_2) $

$= \sum\limits_{i=0}^{k}\langle \textsf{Decomp}^{\beta, l}(\textsf{ct}_i), \bar{\textsf{ct}}_i \rangle $

$\approx \textsf{LWE}_{\vec{s}, \sigma}(\Delta m_1 m_2)$
\end{tcolorbox}

这意味着，将两个 TFHE 密文（一个为 LWE，另一个为 GSW）相乘，并解密所得的 LWE 密文，结果与将两个原始明文相乘相同。


%\textbf{\underline{Proof}}
\begin{myproof}
\begin{enumerate}
\item $\sum\limits_{i=0}^{k}\langle \textsf{Decomp}^{\beta, l}(\textsf{ct}_i), \bar{\textsf{ct}}_i \rangle$ \\
$= \langle \textsf{Decomp}^{\beta, l}(a_0), \bar{\textsf{ct}}_0 \rangle + \langle \textsf{Decomp}^{\beta, l}(a_1), \bar{\textsf{ct}}_1 \rangle + \gap{$\cdots$} + \langle \textsf{Decomp}^{\beta, l}(a_{k-1}), \bar{\textsf{ct}}_{k-1} \rangle + \langle \textsf{Decomp}^{\beta, l}(b), \bar{\textsf{ct}}_k \rangle$ \\
\textcolor{red}{ $\rhd$ 展开两个向量的点积}
\item 对于 $i = k$： \\
$\textsf{Decomp}^{\beta, l}(b) = (b_1, b_2, \ldots , b_l)$，其中 $b = b_1\dfrac{q}{\beta^1} + b_2\dfrac{q}{\beta^2} + \cdots + b_l\dfrac{q}{\beta^l}$ \textcolor{red}{ $\rhd$ 由 \autoref{subsec:poly-decomp} 可得}\\
$\bar{\textsf{ct}}_k = \textsf{Lev}_{\vec{s}, \sigma}^{\beta, l}(m_2) = \left(\textsf{LWE}_{\vec{s}, \sigma}\left(m_{2}\dfrac{q}{\beta^1} + e_{2, 1}\right), \textsf{LWE}_{\vec{s}, \sigma}\left(m_{2}\dfrac{q}{\beta^2} + e_{2, 2}\right), \ldots, \textsf{LWE}_{\vec{s}, \sigma}\left(m_{2}\dfrac{q}{\beta^l} + e_{2, l}\right) \right)$ \\
$ $

因此： \\
$\langle \textsf{Decomp}^{\beta, l}(b), \bar{\textsf{ct}}_k \rangle$\\
$= b_1 \cdot \textsf{LWE}_{\vec{s}, \sigma} \left (m_{2}\dfrac{q}{\beta^1}  + e_{2, 1}\right ) + b_2 \cdot \textsf{LWE}_{\vec{s}, \sigma} \left (m_{2}\dfrac{q}{\beta^2}  + e_{2, 2}\right ) + \gap{$\cdots$} + b_l \cdot \textsf{LWE}_{\vec{s}, \sigma} \left (m_{2}\dfrac{q}{\beta^l}  + e_{2, l}\right)$ \\
$= \textsf{LWE}_{\vec{s}, \sigma} \left (b_1m_{2}\dfrac{q}{\beta^1}  + b_1e_{2, 1}\right ) + \textsf{LWE}_{\vec{s}, \sigma} \left (b_2m_{2}\dfrac{q}{\beta^2}  + b_2e_{2,2}\right ) + \gap{$\cdots$} + \textsf{LWE}_{\vec{s}, \sigma} \left (b_lm_{2}\dfrac{q}{\beta^l} + b_le_{2, l}\right)$ \textcolor{red}{ $\rhd$ 由 \autoref{sec:glwe-mult-plain} 可得} \\
$= \textsf{LWE}_{\vec{s}, \sigma} \left (b_1m_{2}\dfrac{q}{\beta^1} + b_2m_{2}\dfrac{q}{\beta^2} + \gap{$\cdots$} + b_lm_{2}\dfrac{q}{\beta^l} + b_1e_{2,1} + \cdots + b_le_{2,l}\right)$ \textcolor{red}{ $\rhd$ 由 \autoref{sec:glwe-add-cipher} 可得}\\
$= \textsf{LWE}_{\vec{s}, \sigma} \left (m_{2} \cdot \left ( b_1\dfrac{q}{\beta^1} + b_2\dfrac{q}{\beta^2} + \gap{$\cdots$} + b_l\dfrac{q}{\beta^l} + e^{\langle k \rangle} \right)\right)$ \textcolor{red}{ $\rhd$ 其中 $e^{\langle k \rangle} = \sum\limits_{i=1}^{l}b_ie_{2, i}$}\\
$= \textsf{LWE}_{\vec{s}, \sigma} (m_{2}b + e^{\langle k \rangle})$ \textcolor{red}{ $\rhd$ 由 \autoref{subsec:poly-decomp} 可得}

\item 对于 $0 \leq i \leq (k - 1)$： \\
$\textsf{Decomp}^{\beta, l}(a_i) = (a_{\langle i, 1 \rangle}, a_{\langle i, 2 \rangle}, \gap{$\cdots$}, a_{\langle i, l \rangle})$，其中 $a_i = a_{\langle i, 1 \rangle}\dfrac{q}{\beta^1} + a_{\langle i, 2 \rangle}\dfrac{q}{\beta^2} + \gap{$\cdots$} + a_{\langle i, l \rangle}\dfrac{q}{\beta^l}$ \\
$\bar{\textsf{ct}}_i =  \textsf{Lev}_{\vec{s}, \sigma}^{\beta, l}(-s_im_2) $

$= \left(\textsf{LWE}_{\vec{s}, \sigma}\left(-s_im_2\dfrac{q}{\beta^1} + e_{i, 1}\right), \textsf{LWE}_{\vec{s}, \sigma}\left(-s_im_2\dfrac{q}{\beta^2} + e_{i, 2}\right), \gap{$\cdots$}, \textsf{LWE}_{\vec{s}, \sigma}\left(-s_im_2\dfrac{q}{\beta^l} + e_{i, l}\right) \right)$ \\

$ $
因此： \\
$\langle \textsf{Decomp}^{\beta, l}(a_0), \bar{\textsf{ct}}_0 \rangle + \langle \textsf{Decomp}^{\beta, l}(a_1), \bar{\textsf{ct}}_1 \rangle + \gap{$\cdots$} + \langle \textsf{Decomp}^{\beta, l}(a_{k-1}), \bar{\textsf{ct}}_{k-1} \rangle$ \\
$= \sum\limits_{i=0}^{k-1}\langle \textsf{Decomp}^{\beta, l}(a_i), \bar{\textsf{ct}}_i \rangle$ \\
$= \sum\limits_{i=0}^{k-1}\Bigg(a_{\langle i, 1\rangle} \cdot \textsf{LWE}_{\vec{s}, \sigma}\left(-s_im_2\dfrac{q}{\beta^1} + e_{i,1}\right) + a_{\langle i, 2\rangle} \cdot \textsf{LWE}_{\vec{s}, \sigma}\left(-s_im_2\dfrac{q}{\beta^2} + e_{i,2}\right) + $

$ \cdots +  a_{\langle i, l\rangle} \cdot \textsf{LWE}_{\vec{s}, \sigma}\left(-s_im_2\dfrac{q}{\beta^l} + e_{i,l}\right)\Bigg)$ \\
$= \sum\limits_{i=0}^{k-1}\Bigg(\textsf{LWE}_{\vec{s}, \sigma}\left(-a_{\langle i, 1\rangle}s_im_2\dfrac{q}{\beta^1} + a_{\langle i, 1 \rangle}e_{i,1}\right) + \textsf{LWE}_{\vec{s}, \sigma}\left(-a_{\langle i, 2\rangle}s_im_2\dfrac{q}{\beta^2} + a_{\langle i, 2 \rangle}e_{i,2}\right) + $

$\gap{$\cdots$} + \textsf{LWE}_{\vec{s}, \sigma}\left(-a_{\langle i, l\rangle}s_im_2\dfrac{q}{\beta^l} + a_{\langle i, l \rangle}e_{i,l}\right)\Bigg)$ \\
$= \sum\limits_{i=0}^{k-1}\textsf{LWE}_{\vec{s}, \sigma}\left(-a_{\langle i, 1\rangle}s_im_2\dfrac{q}{\beta^1} -a_{\langle i, 2\rangle}s_im_2\dfrac{q}{\beta^2} + \gap{$\cdots$} -a_{\langle i, l\rangle}s_im_2\dfrac{q}{\beta^l} + a_{\langle i, 1 \rangle}e_{i,1} + \cdots + a_{\langle i, l \rangle}e_{i,l}\right)$ \\
$= \sum\limits_{i=0}^{k-1}\textsf{LWE}_{\vec{s}, \sigma}\left(-s_im_2 \cdot \left(a_{\langle i, 1\rangle}\dfrac{q}{\beta^1} + a_{\langle i, 2\rangle}\dfrac{q}{\beta^2} + \gap{$\cdots$} + a_{\langle i, l\rangle}\dfrac{q}{\beta^l}\right) + a_{\langle i, 1 \rangle}e_{i,1} + \cdots + a_{\langle i, l \rangle}e_{i,l}\right)$ \\
$= \sum\limits_{i=0}^{k-1}\textsf{LWE}_{\vec{s}, \sigma}(-s_im_2a_i + e^{\langle i \rangle})$ \textcolor{red}{ $\rhd$ 其中 $e^{\langle i \rangle} = a_{\langle i, 1 \rangle}e_{i,1} + \cdots + a_{\langle i, l \rangle}e_{i,l}$}
\item 根据步骤 2 和 3， \\
$\sum\limits_{i=0}^{k}\langle \textsf{Decomp}^{\beta, l}(\textsf{ct}_i), \bar{\textsf{ct}}_i \rangle$ \\
$= \sum\limits_{i=0}^{k-1}\textsf{LWE}_{\vec{s}, \sigma}(-s_im_2a_i + e^{\langle i \rangle}) + \textsf{LWE}_{\vec{s}, \sigma} (m_{2}b + e^{\langle k \rangle})$ \\ 
$= \textsf{LWE}_{\vec{s}, \sigma}\left(\sum\limits_{i=0}^{k-1}(-s_im_2a_i) + m_{2}b + \sum\limits_{i=0}^ke^{\langle i \rangle}\right)$ \textcolor{red}{ $\rhd$ 两个 LWE 密文相加} \\
$= \textsf{LWE}_{\vec{s}, \sigma}\left(m_2b - \sum\limits_{i=0}^{k-1}m_2a_is_i + \sum\limits_{i=0}^ke^{\langle i \rangle}\right)$  \\ 
$= \textsf{LWE}_{\vec{s}, \sigma}\left(m_2(b - \sum\limits_{i=0}^{k-1}a_is_i) + \sum\limits_{i=0}^ke^{\langle i \rangle}\right)$\\
$= \textsf{LWE}_{\vec{s}, \sigma}\left(m_2(\Delta m_1 + e_1) + \sum\limits_{i=0}^ke^{\langle i \rangle}\right)$ \\
$= \textsf{LWE}_{\vec{s}, \sigma}\left(\Delta m_1m_2 + m_2e_1 + \sum\limits_{i=0}^ke^{\langle i \rangle}\right)$ \\
$\approx \textsf{LWE}_{\vec{s}, \sigma}(\Delta m_1m_2)$ \textcolor{red}{ $\rhd$ 因为 $m_2e_1 + \sum\limits_{i=0}^ke^{\langle i \rangle}$ 很小，所以 $m_2e_1$ 也很小}
\end{enumerate}

\end{myproof}

%This issue of noise growth is similar to that in ciphertext-to-plaintext multiplication (\autoref{subsubsec:glwe-mult-plain-discussion}). 
为了抑制噪声增长，需要进行噪声自举（将在 \autoref{subsec:tfhe-noise-bootstrapping} 中介绍）。

\subsubsection{推广到 GLWE 与 GGSW 的乘法}
\label{subsubsec:tfhe-glwe-to-ggsw-multiplication}

可以将 TFHE 的 LWE 与 GSW 乘法进一步推广为以下两个密文之间的 GLWE 与 GGSW 乘法：$\textsf{GLWE}_{\vec{S}, \sigma}(\Delta M_1) \cdot \textsf{GGSW}_{\vec{S}, \sigma}^{\beta, l}(M_2)$，其中 $M_1$、$M_2$ 和 $S$ 都是 $(n-1)$ 次多项式。


$ $

\noindent 第一项 $\textsf{GLWE}_{\vec{S}, \sigma}(\Delta M_1)$ 可由以下任一种方式得到：
\begin{itemize}
\item 对明文 $M_1$ 进行一次新的 GLWE 加密（\autoref{subsec:glwe-enc}）。
\item 对两个 GLWE 密文执行同态加法所得的结果（\autoref{sec:glwe-add-cipher}）。
\item 对一个 GLWE 密文和一个明文执行同态乘法所得的结果（\autoref{sec:glwe-mult-plain}）。
\end{itemize}

$ $

\noindent 第二项 $\textsf{GGSW}_{\vec{S}, \sigma}^{\beta, l}(M_2)$ 可由以下任一种方式得到：
\begin{itemize}
\item 对明文 $M_2$ 进行一次新的 GGSW 加密（\autoref{subsec:ggsw-enc}）。
\item 通过\textit{电路自举}将 $\textsf{GLWE}_{\vec{S}, \sigma}(\Delta M_2)$ 转换为 $\textsf{GGSW}_{\vec{S}, \sigma}^{\beta, l}(M_2)$（将在未来版本中介绍）。
\end{itemize}

$ $

\noindent 回顾以下关系：

$\textsf{GLWE}_{\vec{S}, \sigma}(\Delta M_1) = (A_0, A_1, \gap{$\cdots$}, A_{k-1}, B) \in \mathcal{R}_{n, q}^{k + 1}$，其中 $B = \sum\limits_{i=0}^{k-1}(A_i \cdot S_i) + \Delta M_1 + E$ \\ \textcolor{red}{ $\rhd$ 由 \autoref{subsec:glwe-enc} 可得}

$ $

$\textsf{GGSW}_{\vec{S}, \sigma}^{\beta, l}(M_2) = \Bigr \{ \{ \textsf{GLev}_{\vec{S}, \sigma}^{\beta, l} (-S_i \cdot M_2)  \}_{i=0}^{k-1}, \textsf{GLev}_{\vec{S}, \sigma}^{\beta, l}(M_2) \Bigl \} \in \mathcal{R}_{\langle n, q \rangle }^{(k+1) \cdot l \cdot (k'+1)}$ \textcolor{red}{ $\rhd$ 由 \autoref{subsec:ggsw-enc} 可得}

$ $

\noindent 使用以下记号：

$\textsf{GGSW}_{\vec{S}, \sigma}^{\beta, l}(M_2) = {\bar{C}} = (\bar{C_0},. \bar{C_1}, \gap{$\cdots$} \bar{C_k})$ 

$\bar{C_i} = \textsf{GLev}_{\vec{S}, \sigma}^{\beta, l}(-S_i \cdot M_2)$，其中 $0 \leq i \leq (k-1)$

$\bar{C_k} = \textsf{GLev}_{\vec{S}, \sigma}^{\beta, l}(M_2)$

$\textsf{ct} = \textsf{GLWE}_{\vec{S}, \sigma}(\Delta M_1) = (C_0, C_1, \gap{$\cdots$}, C_k) = (A_0, A_1, \gap{$\cdots$}, A_{k-1}, B)$

$ $

\noindent 定义以下 TFHE 密文乘法运算：

$\textsf{ct} \cdot {\bar{C}} = \sum\limits_{i=0}^{k}\langle \textsf{Decomp}^{\beta, l}(C_i), \bar{C_i} \rangle$

$ $

\noindent 于是下式成立：

\begin{tcolorbox}[title={\textbf{\tboxlabel{\ref*{subsubsec:tfhe-glwe-to-ggsw-multiplication}} 推广到 GLWE 与 GGSW 的乘法}}]
$\textsf{ct} = \textsf{GLWE}_{\vec{S}, \sigma}(\Delta M_1) = (A_0, A_1, \cdots, A_{k-1}, B)$

$\bar{C} = \textsf{GGSW}_{\vec{S}, \sigma}^{\beta, l}(M_2)$

$ = \bm(\textsf{GLev}_{\vec{S}, \sigma}^{\beta, l}(-S_0\cdot M_2), \textsf{GLev}_{\vec{S}, \sigma}^{\beta, l}(-S_1\cdot M_2), \cdots, \textsf{GLev}_{\vec{S}, \sigma}^{\beta, l}(-S_{k-1}\cdot M_2), \textsf{GLev}_{\vec{S}, \sigma}^{\beta, l}(M_2)\bm)$

$ $

$\textsf{GLWE}_{\vec{S}, \sigma}(\Delta M_1) \cdot \textsf{GGSW}_{\vec{S}, \sigma}^{\beta, l}(M_2) = \sum\limits_{i=0}^{k}\langle \textsf{Decomp}^{\beta, l}(C_i), \bar{C_i} \rangle \approx \textsf{GLWE}_{\vec{S}, \sigma}(\Delta M_1 M_2)$
\end{tcolorbox}

这意味着，将两个 TFHE 密文（一个为 GLWE，另一个为 GGSW）相乘，并解密所得的 GLWE 密文，结果与将两个原始明文相乘相同。


%\textbf{\underline{Proof}}
\begin{myproof}
\begin{enumerate}
\item $\sum\limits_{i=0}^{k}\langle \textsf{Decomp}^{\beta, l}(C_i), \bar{C_i} \rangle$ \\
$= \langle \textsf{Decomp}^{\beta, l}(A_0), \bar{C_0} \rangle + \langle \textsf{Decomp}^{\beta, l}(A_1), \bar{C_1} \rangle + \gap{$\cdots$} + \langle \textsf{Decomp}^{\beta, l}(A_{k-1}), \bar{C}_{k-1} \rangle + \langle \textsf{Decomp}^{\beta, l}(B), \bar{C_k} \rangle$ \\
\textcolor{red}{ $\rhd$ 展开两个向量的点积}
\item 对于 $i = k$： \\
$\textsf{Decomp}^{\beta, l}(B) = (B_1, B_2, \gap{$\cdots$}, B_l)$，其中 $B = B_1\dfrac{q}{\beta^1} + B_2\dfrac{q}{\beta^2} + \gap{$\cdots$} + B_l\dfrac{q}{\beta^l}$ \textcolor{red}{ $\rhd$ 由 \autoref{subsec:poly-decomp} 可得}\\
$\bar{C}_k = \textsf{GLev}_{\vec{S}, \sigma}^{\beta, l}(M_2) = \left(\textsf{GLWE}_{\vec{S}, \sigma}\left(M_{2}\dfrac{q}{\beta^1}\right), \textsf{GLWE}_{\vec{S}, \sigma}\left(M_{2}\dfrac{q}{\beta^2}\right), \gap{$\cdots$}, \textsf{GLWE}_{\vec{S}, \sigma}\left(M_{2}\dfrac{q}{\beta^l}\right) \right)$ 

\textcolor{red}{ $\rhd$ 为简化表达，省略每个 GLWE 密文中的噪声项 $E_{2,1}, \ldots, E_{2,l}$} \\
$ $

因此： \\
$\langle \textsf{Decomp}^{\beta, l}(B), \bar{C_k} \rangle$\\
$= B_1 \cdot \textsf{GLWE}_{\vec{S}, \sigma} \left (M_{2}\dfrac{q}{\beta^1} \right ) + B_2 \cdot \textsf{GLWE}_{\vec{S}, \sigma} \left (M_{2}\dfrac{q}{\beta^2} \right ) + \gap{$\cdots$} + B_l \cdot \textsf{GLWE}_{\vec{S}, \sigma} \left (M_{2}\dfrac{q}{\beta^l}\right)$ \\
$= \textsf{GLWE}_{\vec{S}, \sigma} \left (B_1M_{2}\dfrac{q}{\beta^1} \right ) + \textsf{GLWE}_{\vec{S}, \sigma} \left (B_2M_{2}\dfrac{q}{\beta^2} \right ) + \gap{$\cdots$} + \textsf{GLWE}_{\vec{S}, \sigma} \left (B_lM_{2}\dfrac{q}{\beta^l}\right)$ \textcolor{red}{ $\rhd$ 由 \autoref{sec:glwe-mult-plain} 可得} \\
$= \textsf{GLWE}_{\vec{S}, \sigma} \left (B_1M_{2}\dfrac{q}{\beta^1} + B_2M_{2}\dfrac{q}{\beta^2} + \gap{$\cdots$} + B_lM_{2}\dfrac{q}{\beta^l}\right)$ \textcolor{red}{ $\rhd$ 由 \autoref{sec:glwe-add-cipher} 可得}\\
$= \textsf{GLWE}_{\vec{S}, \sigma} \left (M_{2} \cdot \left ( B_1\dfrac{q}{\beta^1} + B_2\dfrac{q}{\beta^2} + \gap{$\cdots$} + B_l\dfrac{q}{\beta^l} \right)\right)$ \\
$= \textsf{GLWE}_{\vec{S}, \sigma} (M_{2}B)$ \textcolor{red}{ $\rhd$ 由 \autoref{subsec:poly-decomp} 可得}
\item 对于 $0 \leq i \leq (k - 1)$： \\
$\textsf{Decomp}^{\beta, l}(A_i) = (A_{\langle i, 1 \rangle}, A_{\langle i, 2 \rangle}, \gap{$\cdots$}, A_{\langle i, l \rangle})$，其中 $A_i = A_{\langle i, 1 \rangle}\dfrac{q}{\beta^1} + A_{\langle i, 2 \rangle}\dfrac{q}{\beta^2} + \gap{$\cdots$} + A_{\langle i, l \rangle}\dfrac{q}{\beta^l}$ \\
$\bar{C_i} =  \textsf{GLev}_{\vec{S}, \sigma}^{\beta, l}(-S_iM_2) = \left(\textsf{GLWE}_{\vec{S}, \sigma}\left(-S_iM_2\dfrac{q}{\beta^1}\right), \textsf{GLWE}_{\vec{S}, \sigma}\left(-S_iM_2\dfrac{q}{\beta^2}\right), \gap{$\cdots$}, \textsf{GLWE}_{\vec{S}, \sigma}\left(-S_iM_2\dfrac{q}{\beta^l}\right) \right)$ \\

$ $
因此： \\
$\langle \textsf{Decomp}^{\beta, l}(A_0), \bar{C_0} \rangle + \langle \textsf{Decomp}^{\beta, l}(A_1), \bar{C_1} \rangle + \gap{$\cdots$} + \langle \textsf{Decomp}^{\beta, l}(A_{k-1}), \bar{C}_{k-1} \rangle$ \\
$= \sum\limits_{i=0}^{k-1}\langle \textsf{Decomp}^{\beta, l}(A_i), \bar{C_i} \rangle$ \\
$= \sum\limits_{i=0}^{k-1}\left(A_{\langle i, 1\rangle} \cdot \textsf{GLWE}_{\vec{S}, \sigma}\left(-S_iM_2\dfrac{q}{\beta^1}\right) + A_{\langle i, 2\rangle} \cdot \textsf{GLWE}_{\vec{S}, \sigma}\left(-S_iM_2\dfrac{q}{\beta^2}\right) + \gap{$\cdots$} + A_{\langle i, l\rangle} \cdot \textsf{GLWE}_{\vec{S}, \sigma}\left(-S_iM_2\dfrac{q}{\beta^l}\right)\right)$ \\
$= \sum\limits_{i=0}^{k-1}\left(\textsf{GLWE}_{\vec{S}, \sigma}\left(-A_{\langle i, 1\rangle}S_iM_2\dfrac{q}{\beta^1}\right) + \textsf{GLWE}_{\vec{S}, \sigma}\left(-A_{\langle i, 2\rangle}S_iM_2\dfrac{q}{\beta^2}\right) + \gap{$\cdots$} + \textsf{GLWE}_{\vec{S}, \sigma}\left(-A_{\langle i, l\rangle}S_iM_2\dfrac{q}{\beta^l}\right)\right)$ \\
$= \sum\limits_{i=0}^{k-1}\textsf{GLWE}_{\vec{S}, \sigma}\left(-A_{\langle i, 1\rangle}S_iM_2\dfrac{q}{\beta^1} + -A_{\langle i, 2\rangle}S_iM_2\dfrac{q}{\beta^2} + \gap{$\cdots$} + -A_{\langle i, l\rangle}S_iM_2\dfrac{q}{\beta^l}\right)$ \\
$= \sum\limits_{i=0}^{k-1}\textsf{GLWE}_{\vec{S}, \sigma}\left(-S_iM_2 \cdot \left(A_{\langle i, 1\rangle}\dfrac{q}{\beta^1} + A_{\langle i, 2\rangle}\dfrac{q}{\beta^2} + \gap{$\cdots$} + A_{\langle i, l\rangle}\dfrac{q}{\beta^l}\right)\right)$ \\
$= \sum\limits_{i=0}^{k-1}\textsf{GLWE}_{\vec{S}, \sigma}(-S_iM_2A_i)$
\item 根据步骤 2 和 3， \\
$\sum\limits_{i=0}^{k}\langle \textsf{Decomp}^{\beta, l}(C_i), \bar{C_i} \rangle$ \\
$= \sum\limits_{i=0}^{k-1}\textsf{GLWE}_{\vec{S}, \sigma}(-S_iM_2A_i) + \textsf{GLWE}_{\vec{S}, \sigma} (M_{2}B)$ \\ 
$= \textsf{GLWE}_{\vec{S}, \sigma}\Big(\sum\limits_{i=0}^{k-1}(-S_iM_2A_i) + M_{2}B\Big)$ \textcolor{red}{ $\rhd$ 两个 GLWE 密文相加} \\
$= \textsf{GLWE}_{\vec{S}, \sigma}\Big(BM_2 - \sum\limits_{i=0}^{k-1}M_2A_iS_i\Big)$  \\ 
$= \textsf{GLWE}_{\vec{S}, \sigma}\Big(M_2(B - \sum\limits_{i=0}^{k-1}A_iS_i)\Big)$\\
$= \textsf{GLWE}_{\vec{S}, \sigma}(M_2(\Delta M_1 + E))$ \\
$= \textsf{GLWE}_{\vec{S}, \sigma}(\Delta M_1M_2 + M_2E)$ \\
$\approx \textsf{GLWE}_{\vec{S}, \sigma}(\Delta M_1M_2)$ \textcolor{red}{ $\rhd$ 因为 $E$ 很小，所以 $M_2E$ 也很小}
\end{enumerate}

\end{myproof}

\subsection{系数提取}
\label{subsec:tfhe-extraction}

\textbf{- 参考资料：}
\href{https://www.zama.ai/post/tfhe-deep-dive-part-4}{深入理解 TFHE：第四部分——可编程自举}~\cite{tfhe-4}

$ $

在 TFHE 中，系数提取是指从以 GLWE 密文形式加密的多项式中提取某个系数的过程。提取出的系数仍以 LWE 密文的形式表示（\autoref{sec:lwe}）。 %We will explain how to extract the coefficient of a plaintext polynomial $M$ from a GLWE ciphertext and RLWE ciphertext, respectively.


%\subsubsection{Coefficient Extraction from a GLWE Ciphertext}
%\label{subsubsec:tfhe-extraction-glwe}

注意，在 GLWE 密码系统中，明文 $M$ 被编码为一个多项式，其中每个系数都编码一个明文值 $m_0, m_1, \cdots, m_{n-1}$。

%\subsubsection{Overview}
%\label{subsec:tfhe-extraction-overview}

假设 GLWE 密文的设置如下： \\
$M = \sum\limits_{j=0}^{n-1}m_jX^j \in \mathcal{R}_{\langle n, q \rangle}$ 

$S = \left(S_0 = \sum\limits_{j=0}^{n-1}s_{0,j}X^j, S_1 = \sum\limits_{j=0}^{n-1}s_{1,j}X^j, \gap{$\cdots$}, S_{k-1} = \sum\limits_{j=0}^{n-1}s_{k-1,j}X^j \right)$ 

$\textsf{GLWE}_{\vec{S}, \sigma}(\Delta M) = \left(A_0 = \sum\limits_{j=0}^{n-1}a_{0,j}X^j, A_1 = \sum\limits_{j=0}^{n-1}a_{1,j}X^j, \gap{$\cdots$}, A_{k-1} = \sum\limits_{j=0}^{n-1}a_{k-1,j}X^j, B = \sum\limits_{j=0}^{n-1}b_{j}X^j\right)$ 

$B = \sum\limits_{i=0}^{k-1}A_iS_i + \Delta M + E$ 

$E = \sum\limits_{i=0}^{n-1}e_iX^i$ 

$ $

\noindent 注意：

$\Delta M + E = B - \sum\limits_{i=0}^{k-1}A_iS_i$ 

$ = (\Delta m_0 + \Delta m_1X + \gap{$\cdots$} + \Delta m_{n-1}X^{n-1}) + (e_0 + e_1X + \gap{$\cdots$} + e_{n-1}X^{n-1})$

$= (\Delta m_0 + e_0) + (\Delta m_1 + e_1)X + \gap{$\cdots$} + (\Delta m_{n-1} + e_{n-1})X^{n-1}$

$ $

\noindent 该公式还可以写为：

$B - \sum\limits_{i=0}^{k-1}A_iS_i$ 

$ = (b_0 + b_1X + \gap{$\cdots$} + b_{n-1}X^{n-1} )$ 

$ - (a_{0,0} + a_{0,1}X + \gap{$\cdots$} + a_{0, n-1}X^{n-1})(s_{0,0} + s_{0,1}X + \gap{$\cdots$} + s_{0, n-1}X^{n-1})$ 

$ - (a_{1,0} + a_{1,1}X + \gap{$\cdots$} + a_{1, n-1}X^{n-1})(s_{1,0} + s_{1,1}X + \gap{$\cdots$} + s_{1, n-1}X^{n-1})$

$ - \gap{$\cdots$} $ 

$ - (a_{k-1,0} + a_{k-1,1}X + \gap{$\cdots$} + a_{k-1, n-1}X^{n-1})(s_{k-1,0} + s_{k-1,1}X + \gap{$\cdots$} + s_{k-1, n-1}X^{n-1})$ 

$ $

$ = \left(b_0 - \left( \sum\limits_{i=0}^{k-1} \sum\limits_{j=0}^{0}(a_{i,0-j}s_{i,j}) - \sum\limits_{i=0}^{k-1} \sum\limits_{j=1}^{n-1}(a_{i,n-j}s_{i,j}) \right)\right)$

$ + \left(b_1 - \left( \sum\limits_{i=0}^{k-1} \sum\limits_{j=0}^{1}(a_{i,1-j}s_{i,j}) - \sum\limits_{i=0}^{k-1} \sum\limits_{j=2}^{n-1}(a_{i,n+1-j}s_{i,j})   \right) \right)\cdot X$ 



$ + \left(b_2 - \left( \sum\limits_{i=0}^{k-1} \sum\limits_{j=0}^{2}(a_{i,2-j}s_{i,j}) - \sum\limits_{i=0}^{k-1} \sum\limits_{j=3}^{n-1}(a_{i,n+2-j}s_{i,j})   \right) \right)\cdot X^2$ 

$ $

$\gap{$\cdots$}$ 

$ $

$ + \left(b_{n-1} - \left(  \sum\limits_{i=0}^{k-1} \sum\limits_{j=0}^{n-1}(a_{i,n-1-j}s_{i,j}) - \sum\limits_{i=0}^{k-1} \sum\limits_{j=n}^{n-1}(a_{i,n+(n-1)-j}s_{i,j})  \right) \right)\cdot X^{n-1}$ 

\textcolor{red}{ $\rhd$ 按相同指数合并各项}

$ $


$ $

$= \sum\limits_{h=0}^{n-1}  \left(b_h - \left(  \sum\limits_{i=0}^{k-1} \sum\limits_{j=0}^{h}(a_{i,h-j}s_{i,j}) - \sum\limits_{i=0}^{k-1} \sum\limits_{j=h+1}^{n-1}(a_{i,n+h-j}s_{i,j})  \right) \right)\cdot X^{h}  $


$ $

$= \sum\limits_{h=0}^{n-1}  C_h \cdot X^{h}  $，其中 $C_h = b_h - \left(  \sum\limits_{i=0}^{k-1} \sum\limits_{j=0}^{h}(a_{i,h-j}s_{i,j}) - \sum\limits_{i=0}^{k-1} \sum\limits_{j=h+1}^{n-1}(a_{i,n+h-j}s_{i,j})  \right)$

$ $

\noindent 注意，在上述 $(n-1)$ 次多项式中，每个 $X^h$ 项的系数 $C_h$ 都可以表示为如下 LWE 密文 $\textsf{ct}_h$：

$S' = (s_{0,0}, s_{0,1}, \gap{$\cdots$}, s_{0,n-1}, s_{1,0}, s_{1,1}, \gap{$\cdots$}, s_{1, n-1}, \gap{$\cdots$}, s_{k-1, n-1}) = (s'_0, s'_1, \gap{$\cdots$}, s'_{nk-1} ) \in \mathbb{Z}_q^{nk}$


$\textsf{ct}_h = (a'_0, a'_1, \gap{$\cdots$}, a'_{nk-1}, b_h) \in \mathbb{Z}_q^{nk + 1}$


\[
    \text{，其中 } a'_{n \cdot i + j} =   
\begin{cases}
    a_{i,h - j} \text{（若 } 0 \leq j \leq h\text{）}\\
    -a_{i,n + h - j} \text{（若 } h+1 \leq j \leq n-1\text{）}\\
\end{cases}
\centering , b_h \text{ 直接取自多项式 } B
\]


\noindent 注意 $b_h - \sum\limits_{i=0}^{nk-1}s'_ia'_i = \Delta m_h + e_h$。这意味着可以用 $C_h$ 的加密形式 $\textsf{LWE}_{\vec{s}_{'}, \sigma}(\Delta m_h)$ 来替代它，即用加密了 $M$ 的第 $h$ 个系数的 LWE 密文 $\textsf{ct}_h$ 来表示。因此，我们从 $\textsf{GLWE}_{\vec{S}, \sigma}(\Delta M)$ 中提取出了 $\textsf{LWE}_{\vec{s}_{'}, \sigma}(\Delta m_h)$。这一操作称为系数提取；它只是重新排列 GLWE 密文多项式的系数以提取一个 LWE 密文，因此不会增加任何噪声。

得到 $\textsf{LWE}_{\vec{s}_{'}, \sigma}(\Delta m_h)$ 后，即可执行 $\vec{s}_{'} \rightarrow \vec{s}$ 的密钥切换（\autoref{subsec:tfhe-key-switching}）。

\begin{tcolorbox}[title={\textbf{\tboxlabel{\ref*{subsec:tfhe-extraction}} GLWE 密文的系数提取}}]
给定以下 GLWE 密文：

$M = \sum\limits_{j=0}^{n-1}m_jX^j \in \mathcal{R}_{\langle n, t \rangle}$

$\vec{S} = \left(S_0 = \sum\limits_{j=0}^{n-1}s_{0,j}X^j, S_1 = \sum\limits_{j=0}^{n-1}s_{1,j}X^j, \gap{$\cdots$}, S_{k-1} = \sum\limits_{j=0}^{n-1}s_{k-1,j}X^j \right)$

$\textsf{GLWE}_{\vec{S}, \sigma}(\Delta M) = \left(A_0 = \sum\limits_{j=0}^{n-1}a_{0,j}X^j, A_1 = \sum\limits_{j=0}^{n-1}a_{1,j}X^j, \gap{$\cdots$}, A_{k-1} = \sum\limits_{j=0}^{n-1}a_{k-1,j}X^j, B = \sum\limits_{j=0}^{n-1}b_{j}X^j\right)$

$B = \sum\limits_{i=0}^{k-1}A_iS_i + \Delta M + E \bmod q$, \text{ } $E = \sum\limits_{i=0}^{n-1}e_iX^i$

$ $

$\textsf{LWE}_{\vec{s}_{'}, \sigma}(\Delta m_h)$ 是一个 LWE 密文，它加密了 $\Delta M$ 的第 $h$ 个系数（即 $\Delta m_h$）。可以按如下方式从 $\textsf{GLWE}_{\vec{S}, \sigma}(\Delta M)$ 中提取 $\textsf{LWE}_{\vec{s}_{'}, \sigma}(\Delta m_h)$：

$ $

$\vec{s}_{'} = (s_{0,0}, s_{0,1}, \gap{$\cdots$}, s_{0,n-1}, s_{1,0}, s_{1,1}, \gap{$\cdots$}, s_{1, n-1}, \gap{$\cdots$}, s_{k-1, n-1}) = (s'_0, s'_1, \gap{$\cdots$}, s'_{nk-1} ) \in \mathbb{Z}_q^{nk}$

$\textsf{LWE}_{\vec{s}_{'}, \sigma}(\Delta m_h) = (a_0', a_1', \gap{$\cdots$} , a_{nk-1}', b_h) \in \mathbb{Z}_q^{nk + 1}$


\[
    \text{，其中 } a'_{n \cdot i + j} =   
\begin{cases}
    a_{i,h - j} \text{（若 } 0 \leq j \leq h\text{）}\\
    -a_{i,n + h - j} \text{（若 } h+1 \leq j \leq n-1\text{）}\\
\end{cases}
, b_h \text{ 取自多项式 } B
\]

得到 $\textsf{LWE}_{\vec{s}_{'}, \sigma}(\Delta m_h)$ 后，执行 $\vec{s}_{'} \rightarrow \vec{s}$ 的密钥切换（\autoref{subsec:tfhe-key-switching}）。

\end{tcolorbox}



\subsection{噪声自举}
\label{subsec:tfhe-noise-bootstrapping}


\textbf{- 参考资料：}
\href{https://www.zama.ai/post/tfhe-deep-dive-part-4}{深入理解 TFHE：第四部分——可编程自举}~\cite{tfhe-4}

$ $

连续对 TFHE 密文进行同态加法不一定会使噪声 $e$ 增大，因为 $e$ 是按照高斯分布随机生成的，许多噪声相加后的均值为 0。另一方面，连续进行同态乘法会使噪声增长，因为噪声项也会相乘，从而增大其幅度。因此，需要在噪声侵入存放明文 $m$ 的高位之前，以某种方式将噪声\textit{重置}（即防止 \autoref{fig:scaling} 中红色噪声位溢出到蓝色明文位）。将噪声重新初始化为较小值的过程称为噪声自举。

如本节开头所述，TFHE 使用 LWE（即多项式次数为 0 的 GLWE）对明文进行加密与解密。也就是说，每个明文都是 $m$（一个数），被编码为零次多项式。此外，加密各个 $m$ 的秘密密钥 S 是向量 $ \{s_0, s_1, \text{ } \cdots \text{ }, s_{k-1} \}$，而不是一个多项式。另一方面，TFHE 的噪声自举使用 GLWE 密文之间的同态加法，以及 GLWE 与 GGSW 密文之间的同态乘法。

%Nonetheless, the reason why we learned TFHE ciphertext addition (\autoref{subsec:glwe-add} and \autoref{subsec:glwe-add-plain}), multiplication (\autoref{subsec:tfhe-mult-cipher} and \autoref{subsec:glwe-mult-plain}), and polynomial rotation (\autoref{subsec:coeff-rotation}) in terms of GLWE is because the bootstrapping procedure we will explain in this subsection requires homomorphic addition and multiplication of GLWE ciphertexts.

假设有如下 TFHE 密文：

$ $

$\textsf{LWE}_{\vec{s}, \sigma}(\Delta m) = (a_0, a_1, \gap{$\cdots$} a_{k-1}, b)$

$b = \sum\limits_{i=0}^{k-1} a_is_i + \Delta m + e_b$

$\vec{s} = (s_0, s_1, \gap{$\cdots$} s_{k-1})$

$ $

，其中 $e_b$ 是经过许多次密文（或明文）乘法累积而成的较大噪声。噪声自举的目标是将 $(a_0, a_1, \gap{$\cdots$} a_{k-1}, b)$ 转换为 $(a_0', a_1', \gap{$\cdots$} a_{k-1}', b')$，使其满足：

$ $

$b' = \sum\limits_{i=0}^{k-1} a_i's_i + \Delta m + e_s$

$ $

，其中 $e_s$ 是重新初始化后的噪声。

$ $
\subsubsection{概述}
\label{subsec:bootstrapping-overview}

为了实现噪声自举，需要构造一个专门设计的 $(n-1)$ 次多项式 $V(X)$，称为查找表（Lookup Table，LUT）。在介绍 $V(X)$ 之前，先通过一个初步的 LUT 多项式 $V_q(X)$ 说明其基本思想。设想多项式 $V_q(X)$ 中每个项 $X^{j}$ 的指数 $j = \Delta m_i + e_*$ 是带有噪声 $e_* \in \mathbb{Z}_{\Delta}$ 的明文 $m_i$，而对应的系数 $v_{j} = m_i$ 则是不含噪声的明文。因此，将 $(q-1)$ 次多项式 $V_q(X)$ 定义如下：

$V_q(X) = v_0 + v_1X^1 + v_2X^2 + \gap{$\cdots$} + v_{q-1}X^{q-1}$

\text{ } $= m_0X^{\Delta m_0 + e_0} + m_0X^{\Delta m_0 + e_1} + m_0X^{\Delta m_0 + e_2} + \gap{$\cdots$} + m_0X^{\Delta m_0 + e_{\Delta - 1}}$ \textcolor{red}{ $\rhd$ 共 $\Delta$ 项}

\text{ } $ + \text{ } m_1X^{\Delta m_1 + e_0} + m_1X^{\Delta m_1 + e_1} + m_1X^{\Delta m_1 + e_2} + \gap{$\cdots$} + m_1X^{\Delta m_1 + e_{\Delta - 1}}$ \textcolor{red}{ $\rhd$ 共 $\Delta$ 项}

\text{ } $ + \gap{$\cdots$} $

\text{ } $ + \text{ } m_{t - 1}X^{\Delta m_{t - 1} + e_0} + m_{t - 1}X^{\Delta m_{t - 1} + e_1} + m_{t - 1}X^{\Delta m_{t - 1} + e_2} + \gap{$\cdots$} + m_{t - 1}X^{\Delta m_{t - 1} + e_{\Delta - 1}}$ \textcolor{red}{ $\rhd$ 共 $\Delta$ 项}

$ $

在上式中，$m_i$ 和 $e_k$ 分别遍历所有可能的明文消息和误差值（其中 $m_i \in \mathbb{Z}_t$ 且 $e_k \in \mathbb{Z}_{\Delta}$）。 %Note that $\mathbb{Z}_{q} = \{0, 1, \cdots, q-1\} = \{\Delta m_0 + e_0, \Delta m_0 + e_1, \cdots, \Delta m_{t-1} + e_{\Delta - 1}\}$.

我们将 $V_q(X)$ 设计为具有如下性质：每个项 $v_{j}X^{j}$ 都表示特殊的映射（指数，系数）$= (\Delta m_i + e_{*}, m_i)$，其中 $e_*$ 可以取 $\mathbb{Z}_{\Delta}$ 中的任意值。在 TFHE 的设置阶段，使用新定义的 GLWE 密钥 $\vec{S}_{bk}$ 对 $V_q(X)$ 进行 GLWE 加密；这一\textit{自举密钥}与 LWE 秘密密钥 $\vec{s}$ 不同。$\vec{S}_{bk}$ 是由系数为二进制值的 $(n-1)$ 次多项式构成的列表。随后，在噪声自举阶段，利用多项式系数旋转方法 1（\autoref{subsec:coeff-rotation} 中的小结~\ref*{subsec:coeff-rotation}.1），计算 $V \cdot X^{-(\Delta m + e)} = V'$，将 $V$ 的系数向左旋转 $\Delta m + e$ 个位置。接着，利用系数提取技术（\autoref{subsec:tfhe-extraction}）提取该多项式常数项的系数，即旋转后的多项式 $V'$ 中最左侧零次项的系数。此外，$V_q(X)$ 在 TFHE 设置阶段已被加密为 GLWE 密文，因此从旋转后的 $V'_q(X)$ 中提取出的常数项系数，是对 $m$ 的 LWE 加密（即 $\textsf{LWE}_{\vec{s}, \sigma}(\Delta m)$），其噪声已重新初始化（即充分降低）。

因此，至少在当前阶段，可以从概念上将噪声自举过程理解为：

$ $

\begin{enumerate}
\item \textbf{\underline{输入}：} 加密 $m$ 的带噪声密文 $\textsf{LWE}_{\vec{s}, \sigma}(\Delta m + e)$
\item 将输入转换为 $X^{-(\Delta m + e)}$ 的形式，作为 $V_q(X)$（查找表）的旋转因子。
\item 计算 $V_q \cdot X^{-(\Delta m + e)} = V'_q$，将 $V_q$ 向左旋转 $\Delta m + e$ 个位置。
\item 以 LWE 密文形式提取旋转后的 $V'_q(X)$ 的常数项系数 $m$，即得到 $\textsf{LWE}_{\vec{s}, \sigma}(\Delta m)$。
\item \textbf{\underline{输出}：} 明文 $m$ 的 LWE 加密 $\textsf{LWE}_{\vec{s}, \sigma}(\Delta m)$，其噪声已重新初始化
\end{enumerate}

$ $

输出 $\textsf{LWE}_{\vec{s}, \sigma}(\Delta m)$ 与输入密文加密相同的明文消息，但噪声已充分降低。因此，可以将输出 $\textsf{LWE}_{\vec{s}, \sigma}(\Delta m)$ 用于后续的 TFHE 同态运算，例如加法或乘法。在该噪声自举过程中，多项式 $V_q$ 充当\textit{字典}，保存从带噪声明文 $\Delta m + e$（即以 $\Delta m + e = j$ 表示，其中 $v_jX^j$）到无噪声明文 $m$（即以 $m = v_j$ 表示，其中 $v_jX^j$）的映射。因此，$V_q$ 被称为查找表（LUT）。

那么，$V_q(X)$ 的次数应取多大？为了让 $V_q(X)$ 编码从 $\Delta m + e \in \mathbb{Z}_q$ 到 $m \in \mathbb{Z}_t$ 的所有可能映射，$V_q(X)$ 应为 $(q-1)$ 次多项式。然而，$q$ 非常大，处理一个 $(q-1)$ 次多项式在计算上不可行。因此，实际使用的是一个小得多的多项式 $V(X)$，其次数只有 $n-1$。回顾 TFHE 的设置，有 $n \ll q$。因此，需要以\textit{比例正确}的方式，将较大的密文空间 $\Delta m + e \in \mathbb{Z}_q$ \textit{压缩}到小得多的空间 $\mathbb{Z}_n$，并将压缩后的值编码为 $X^j$ 的指数。为实现这种 $\mathbb{Z}_q \rightarrow \mathbb{Z}_n$ 的按比例压缩，将使用 \autoref{subsec:modulus-switch-lwe} 中介绍的 LWE 模数切换技术。

\subsubsection{噪声自举中的模数切换}
\label{subsec:bootstrapping-modulus-switch}

为了避免使用庞大的 $(q-1)$ 次多项式 $V_q$，我们将 $q$ 个可能的密文元素 $\Delta m + e \in \mathbb{Z}_q$ 压缩为 $(n-1)$ 次多项式 $V$ 的 $n$ 个不同指数；其中，$V$ 中的每个项 $v_jX^j$ 都表示映射 $j \rightarrow v_j$，即从带噪声明文映射到无噪声明文。但是，将 $(n-1)$ 次多项式 $V$ 的系数向左旋转时，$v_jX^j$ 一旦越过 $X^0$ 次项与 $X^{n-1}$ 次项之间的边界，$v_j$ 的符号就会翻转为 $-v_j$（见 \autoref{subsec:coeff-rotation-ex} 中的例子）。由于这种系数符号翻转，$(n-1)$ 次多项式 $V$ 理论上可以编码共 $2n$ 种不同的系数状态：$(v_0, v_1, v_2, \gap{$\cdots$}, v_{n-1}, -v_0, -v_1, \gap{$\cdots$}, -v_{n-1})$。为了将这 $2n$ 个不同的系数中的每一个移到 $V$ 的常数项位置（即将系数 $v_j$ 移到 $V$ 的最左侧项），旋转运算 $V \cdot X^{-j}$ 可以使用 $2n$ 个不同的 $j$ 值，即 $\{0, 1, 2, \gap{$\cdots$}, n-1, n, \gap{$\cdots$}, 2n-1\}$，从而将 $(v_0, v_1, v_2, \gap{$\cdots$}, v_{n-1}, -v_0, -v_1, \gap{$\cdots$}, -v_{n-1})$ 中的每一个系数移到常数项位置。这说明 $V\cdot X^{-j}$ 中的指数 $j$ 可以取 $2n$ 个不同值中的任意一个，以覆盖 $V$ 的全部 $2n$ 种可能的（包括符号翻转后的）系数状态。此外，记住 $j = \Delta m + e$。因此，将对 $\Delta m + e$ 执行 $q \rightarrow 2n$ 的模数切换。利用 LWE 模数切换技术（\autoref{subsec:modulus-switch-lwe}），原始 LWE 密文 $\textsf{LWE}_{\vec{s}, \sigma}(\Delta m + e) = ({a}_0, {a}_1, \gap{$\cdots$} {a}_{k-1}, {b}) \in \mathbb{Z}_q^{k+1}$，即噪声自举过程的初始输入，将转换为：

$ $

$\textsf{LWE}_{\vec{s}, \sigma}(\hat{\Delta} m) = (\hat{a}_0, \hat{a}_1, \gap{$\cdots$} \hat{a}_{k-1}, \hat{b}) \in \mathbb{Z}_{2n}^{k+1}$

$\vec{s} = (s_0, s_1, \gap{$\cdots$} s_{k-1}) \in \mathbb{Z}_2^{k}$  \textcolor{red}{ $\rhd$ 因为每个 $s_i$ 都是二进制值，所以秘密密钥保持不变}

$\hat{\Delta} = \Delta\dfrac{2n}{q} = \dfrac{2n}{t} \in \mathbb{Z}_{2n}$

$\hat{a}_i = \left \lceil a_i\dfrac{2n}{q} \right \rfloor \in \mathbb{Z}_{2n}$

$\hat{e} = \left \lceil e\dfrac{2n}{q} \right \rfloor \in \mathbb{Z}_{2n}$

$\hat{b} = \left \lceil b\dfrac{2n}{q} \right \rfloor \approx \sum\limits_{i=0}^{k-1} \hat{a}_is_i + \hat{\Delta} m + \hat{e}  \in \mathbb{Z}_{2n}$



$ $

\para{$V(X)$ 的次数与安全性}：如果目标只是设计次数最低的多项式 $V$，使其系数能够映射明文 $m$ 的所有可能取值，那么构造一个 ${t}$ 次多项式 $V$ 就足够了。尽管如此，将 $V$ 的次数选为 $2n$ 而不是 ${t}$，是为了保证足够的安全级别：多项式次数 $n$ 越高，方案对攻击的抵抗能力就越强。



\subsubsection{将可用明文空间减半}
\label{subsec:tfhe-zero-padding} 

一个问题是，LUT 多项式 $V(X)$ 的旋转具有\textit{负循环}性质，即 $V(X)\cdot X^{n} = -V(X)$（也就是说，系数每旋转 $n$ 个位置就会翻转符号）。更一般地：


$V(X)\cdot X^{-j}= V(X)\cdot X^{2n - j} = V(X)\cdot X^{4n - j} = \cdots$


$
= V(X)\cdot X^{-(j \bmod 2n)} =  
\begin{cases}
    v_j  + v_{j+1}X + \cdots, \text{当 } 0\leq i<n\\
    -v_{j} - v_{j-1}X - \cdots, \text{当 } n\leq j<2n
\end{cases}
$

$ $

\noindent ，其中 $v_j$ 表示将多项式 $V(X)$ 向左旋转 $j$ 个位置后得到的常数项系数。问题在于，每当旋转越过 $X^0$ 与 $X^{n-1}$ 之间的边界时，$v_j$ 的符号就会翻转。给定经过模数切换的值 $\hat{a}_j$、$\hat{e}$ 和 $\hat{b}$，构造如下 LUT 多项式 $V(X)$：

$ $

$V(X) = v_0 + v_1X^1 + v_2X^2 + \gap{$\cdots$} + v_{n-1}X^{n-1}$

\text{ } $= m_0X^{\hat{\Delta} m_0 + \hat{e}_0} + m_0X^{\hat{\Delta} m_0 + \hat{e}_1} + m_0X^{\hat{\Delta} m_0 + \hat{e}_2} + \gap{$\cdots$} +  m_0X^{\hat{\Delta} m_0 + \hat{e}_{\hat{\Delta} - 1}}$ \textcolor{red}{ $\rhd$ 共 $\hat{\Delta}$ 项}

\text{ } $ + \text{ } m_1X^{\hat{\Delta} m_1 + \hat{e}_0} + m_1X^{\hat{\Delta} m_1 + \hat{e}_1} + m_1X^{\hat{\Delta} m_1 + \hat{e}_2} + \gap{$\cdots$} + m_1X^{\hat{\Delta} m_1 + \hat{e}_{\hat{\Delta} - 1}}$ \textcolor{red}{ $\rhd$ 共 $\hat{\Delta}$ 项}

\text{ } $ + \gap{$\cdots$} $

\text{ } $ + \text{ } m_{t/2 - 1}X^{\hat{\Delta} m_{t/2 - 1} + \hat{e}_0} + m_{t/2 - 1}X^{\hat{\Delta} m_{t/2 - 1} + \hat{e}_1} + m_{t/2 - 1}X^{\hat{\Delta} m_{t/2 - 1} + \hat{e}_2} + \gap{$\cdots$} + m_{t/2 - 1}X^{\hat{\Delta} m_{t/2 - 1} + \hat{e}_{\hat{\Delta} - 1}}$ \textcolor{red}{ $\rhd$ 共 $\hat{\Delta}$ 项}

$ $


回顾，当 $j= \{0, 1, \cdots , n-1\}$ 时，计算 $V(X)\cdot X^{-j}$ 可以将 $V(X)$ 的系数向左旋转 $\{0, 1, \cdots n-1 \}$ 个位置。对 $V(X)$ 执行每一次 $j$ 个位置的旋转，其中 $j= \{0, 1, \cdots , n-1\}$，旋转后的多项式 $V'(X)$ 的常数项系数依次为：

$\overbrace{\overbrace{\underbrace{\Delta m_0, \Delta m_0, \gap{$\cdots$}}_{\text{系数对应项为 } X^{\hat{\Delta}m_0 + \hat{e}_*}}}^{\hat{\Delta} \text{ 次重复}} \overbrace{\underbrace{\Delta m_1, \Delta m_1, \gap{$\cdots$}}_{\text{系数对应项为 } X^{\hat{\Delta}m_1 + \hat{e}_*}}}^{\hat{\Delta} \text{ 次重复}} \gap{$\cdots$} \overbrace{\underbrace{\Delta m_{{t}/{2}-1}, \Delta m_{{t}/{2}-1}, \gap{$\cdots$}}_{\text{系数对应项为 } X^{\hat{\Delta}m_{{t}/{2}-1}+ \hat{e}_*} }}^{\hat{\Delta} \text{ 次重复}}}^{\text{旋转 $j = \{0, 1, \cdots n-1\}$ 个位置后 $V'(X)$ 常数项的系数}}$ 



$ $

在上式中，噪声 $\hat{e}_*$ 的取值范围为 $[0, \hat{\Delta})$。注意，我们将所有指数 $\hat\Delta m_i + \hat{e}_0, \hat\Delta m_i + \hat{e}_1, \cdots,\hat\Delta m_i + \hat{e}_{\hat\Delta - 1}$ 设计为映射到同一个系数值 $m_i$；这是因为，在对应噪声 $\hat{e}_*$ 被消除后，它们解密得到的底层明文 $m_i$ 相同。因此，每个 $m_i$ 都在 $V(X)$ 的系数中连续重复使用 $\hat\Delta$ 次。可以将这种连续重复理解为一种鲁棒性：对于任意噪声 $\hat{e}_* \in \mathbb{Z}_{\hat{\Delta}}$，都能将相应的 $\hat{\Delta}m_i + \hat{e}_*$ 映射到 $m_i$。

到目前为止，上述 $m_0, m_1, \cdots, m_{{t}/{2} - 1}$ 个系数组成的序列，是我们期望 $V'(X)$（即旋转后的多项式）在执行各次 $0, 1, \cdots, n-1$ 旋转时返回的常数项系数（其中每个 $m_i + 1 = m_{i+1}$）。然而，当旋转次数处于 $[n, 2n -1)$ 范围时，系数映射的正确性就会被破坏，因为多项式环的负循环性质使得系数在从 $X^0$ 次项越过边界到达 $X^{n-1}$ 次项时发生符号翻转。具体而言，每旋转 $n, n+1, \cdots, 2n-1$ 个位置后，多项式 $V'(X)$ 的常数项系数如下：

$ $

$\overbrace{\overbrace{\underbrace{-m_{0}, -m_{0}, \gap{$\cdots$}}_{-\text{系数对应项为 } X^{\hat{\Delta}m_{0} + \hat{e}_*}}}^{\hat{\Delta} \text{ 次重复}} \overbrace{\underbrace{-m_{1}, -m_{1}, \gap{$\cdots$}}_{-\text{系数对应项为 } X^{\hat{\Delta}m_{1} + \hat{e}_*}}}^{\hat{\Delta} \text{ 次重复}} \gap{$\cdots$} \overbrace{\underbrace{-m_{t/2 - 1}, -m_{t/2 - 1}, \gap{$\cdots$}}_{-\text{系数对应项为 } X^{\hat{\Delta}m_{t/2 - 1} + \hat{e}_*} }}^{\hat{\Delta} \text{ 次重复}}}^{\text{分别旋转 $n, n+1, \cdots 2n-1$ 个位置后 $V'(X)$ 常数项的系数}}$ 

$ $

由上式可见，旋转后的 $V'(X)$ 的常数项系数呈现旋转周期为 $n$ 的负循环模式：第二组 $n$ 次旋转得到的系数，恰好是第一组 $n$ 次旋转所得系数的相反数。下面说明这种负循环行为为什么会破坏（指数，系数）= $(\hat{\Delta} m + \hat{e}, m)$ 的映射。由于 TFHE 的明文和密文取值均定义在环上，旋转 LUT 多项式 $V(X)$ 时，理想情况下，旋转后的多项式 $V'(X)$ 的常数项值（即映射后的明文值）应按循环方式回绕，呈现环状模式；每个值连续重复 $\hat{\Delta}$ 次，以抵抗最大为 $\hat{e}_* \in Z_{\hat{\Delta}}$ 的噪声。然而，多项式环的负循环性质使后半组旋转的常数项值出现问题，因为它们恰好是前半组对应值的相反数，从而破坏了第一组与第二组数值之间按环回绕的模式。
%Specifically, as we gradually rotate $V$ with the rotation count $[0, 2n)$, the constant term's coefficient (i.e., $m_i$) should be ideally a value that is either the same as before or greater than the previous value ($m_{i-1}$), all the way to $2n$ rotations (to represent a ring pattern). This constraint holds for the first-half $i = [0, n)$, as the rotated $V'(X)$'s constant term's coefficient is $m_0, \cdots, m_1, \cdots, m_{t/2 - 1}$. However, for the second-half $i = [n, 2n)$, the rotated $V'(X)$'s constant term's coefficient is $-m_0, \cdots, -m_1, \cdots, -m_{t/2 - 1}$, which violates our constraint that $m_i \leq m_{i+1}$. Therefore, the combination of the first-half and the second-half domains of $i$ cannot mathematically implement a continuous modulo value range $\mathbb{Z}_t$. 

因此，$V(X)$ 作为完美 LUT 存在局限：它只能在 $i \in \mathbb{Z}_t$ 的一半范围内保持（指数，系数）= $(\hat{\Delta} m + \hat{e}, m)$ 的正确映射，而无法同时覆盖另一半。

\para{解决方法：}
我们已经看到，当 $V(X)$ 的旋转不超过 $n-1$ 个位置（即第一组旋转）时，$V(X)$ 的（指数，系数）= $(\hat{\Delta} m + \hat e, \Delta m)$ 映射仍保持环状一致性。因此，避免 LUT 多项式旋转的负循环问题，最简单的方法是在 TFHE 应用的\textit{设计中}限制噪声自举期间 $V(X)$ 的旋转不超过 $n-1$ 个位置。为此，当 TFHE 应用的计算流程处理明文值时（即在考虑同态运算之前的原始明文计算逻辑中），应确保每个计算步骤的可能输入和输出，仅涉及 $\mathbb{Z}_t$ 内预先指定的连续 $\dfrac{t}{2}$ 个模剩余值。该约束使 $\hat\Delta m + \hat e$ 的可能取值被限制在 $\mathbb{Z}_{2n}$ 内连续的 $n$ 个值中。由于 LUT 多项式 $V(X)$ 通过计算 $V(X)\cdot X^{-(\hat\Delta m + \hat e)}$ 进行旋转，应用设计将 $\hat\Delta m + \hat e$ 限制为至多 $n-1$（在 $2n - 1$ 范围内），便使 $V(X)$ 在噪声自举期间最多旋转 $n-1$ 个位置。这样就能避免出现使系数符号翻转的、有问题的 $\{n, n+1, \gap{$\cdots$}, 2n -1 \}$ 次旋转。

因此，以将应用可用明文值减半、限制为 $\mathbb{Z}_t$ 内连续的 $\dfrac{t}{2}$ 个值为代价，可以避免 $V(X)$ 的负循环旋转问题，从而保持 $V(X)$ 的（指数，系数）= $(\hat{\Delta} m + \hat e, m)$ 映射正确。

\begin{comment}
与此同时，密文空间 $q$ 保持不变。例如，取简单情形 $t = 8$，TFHE 应用可以设计如下编码方案：

%Note that during the blind rotation (\autoref{subsec:bootstrapping-blind-rotation}), each step's cumulative rotation positions of the homomorphic MUX gate may end up rotating $V(X)$ more than $n$ positions. However, enforcing $0 \leq \hat\Delta m + e < n$ guarantees the rotation-finished $V'(X)$ to have the cumulative rotation count to be some integer within $[0, n) \bmod 2n$.

\begin{multicols}{2}
\begin{itemize}
 \item $0000_2 \rightarrow 000_2 (= 0)$
 \item $0001_2 \rightarrow 001_2 (= 1)$
 \item $0010_2 \rightarrow 010_2 (= 2)$
 \item $0011_2 \rightarrow 011_2 (= 3)$
 \item $0100_2 \rightarrow 100_2 (= -4)$
 \item $0101_2 \rightarrow 101_2 (= -3)$
 \item $0110_2 \rightarrow 110_2 (= 2)$
 \item $0111_2 \rightarrow 111_2 (= 1)$
\end{itemize}
\end{multicols}

通过这种二进制编码，可以在 $t$ 个明文状态中选取 $\dfrac{t}{2}$ 个状态来编码有符号整数，此时最高有效位（MSB）始终为 0。
\end{comment}

综合上述考虑，最终的 LUT 多项式 $V(X)$ 如下：

\begin{tcolorbox}[title={\textbf{\tboxlabel{\ref*{subsec:tfhe-zero-padding}} 查找表多项式 $\bm{V(X)}$ 的结构}}]

$ $

$V(X) = v_0 + v_1X^1 + v_2X^2 + \gap{$\cdots$} + v_{n-1}X^{n-1}$

\text{ } $= m_0X^{\hat{\Delta} m_0 + \hat{e}_0} + m_0X^{\hat{\Delta} m_0 + \hat{e}_1} + m_0X^{\hat{\Delta} m_0 + \hat{e}_2} + \gap{$\cdots$} + m_0X^{\hat{\Delta} m_0 + \hat{e}_{\hat{\Delta} - 1}}$

\text{ } $ + \text{ } m_1X^{\hat{\Delta} m_1 + \hat{e}_0} + m_1X^{\hat{\Delta} m_1 + \hat{e}_1} + m_1X^{\hat{\Delta} m_1 + \hat{e}_2} + \gap{$\cdots$} + m_1X^{\hat{\Delta} m_1 + \hat{e}_{\hat{\Delta} - 1}}$

\text{ } $ + \gap{$\cdots$} $

\text{ } $ + m_{{t}/{2} - 1}X^{\hat{\Delta} m_{{t}/{2} - 1} + \hat{e}_0} + m_{{t}/{2} - 1}X^{\hat{\Delta} m_{{t}/{2} - 1} + \hat{e}_1} + m_{{t}/{2} - 1}X^{\hat{\Delta} m_{{t}/{2} - 1} + \hat{e}_2}$

\text{ } \text{ } $ + \gap{$\cdots$} + m_{{t}/{2} - 1}X^{\hat{\Delta} m_{{t}/{2} - 1} + \hat{e}_{\hat{\Delta} - 1}}$

$ $

\noindent ，其中 $\hat\Delta = \Delta \cdot \dfrac{2n}{q} = \dfrac{2n}{t}$。应用应确保 $m_0, m_1, \cdots, m_{t/2 - 1}$ 是 $\mathbb{Z}_t$ 内一些连续的模 $\dfrac{t}{2}$ 剩余值。该约束使 $V(X)$ 的旋转位置 $\hat{\Delta} m_i + \hat {e}_{*}$（其中 $e_* \in \mathbb{Z}_{\hat{\Delta}}$）至多有 $n$ 种连续的可能值，从而防止 $V(X)$ 旋转超过一个完整周期并触发负循环问题。


\end{tcolorbox}

LUT 多项式 $V(X)$ 也称为累加器。

\subsubsection{盲旋转}
\label{subsec:bootstrapping-blind-rotation} 

盲旋转是指旋转一个\textit{加密}多项式的系数，使观察者无法得知这些系数旋转了多少个位置。旋转后，也无法看出哪个系数移到了哪个次数的项。盲旋转使用基本的多项式旋转方法 1（\autoref{subsec:coeff-rotation} 中的小结~\ref*{subsec:coeff-rotation}），区别在于计算 $V(X)\cdot X^{-i}$ 是以同态方式完成的。

注意 $X^{\hat{\Delta} m + \hat{e}_b} = X^{\hat{b} - \sum_{i=0}^{k-1}{\hat{a}_is_i}}$，因此可以通过计算 $V \cdot X^{-(\hat{b} - \sum_{i=0}^{k-1}{\hat{a}_is_i})} = V \cdot X^{-\hat{b} + \sum_{i=0}^{k-1}{\hat{a}_is_i}}$ 来旋转 $V$。实际上，在不知道 LWE 秘密密钥 $S$ 的情况下，无法直接计算 LWE 解密公式 $-\hat{b} + \sum_{i=0}^{k-1}{\hat{a}_is_i}$（或 $\hat{b} - \sum_{i=0}^{k-1}{\hat{a}_is_i}$）。不过，只要在 TFHE 设置阶段给定 $\{GGSW_{\vec{S}_{bk}, \sigma}^{\beta, l}(s_i)\}_{i=0}^{k-1}$，便存在一种数学方法，使我们无须知道秘密密钥 $S$ 就能计算 $V \cdot X^{-\hat{b} + \sum_{i=0}^{k-1}{\hat{a}_is_i}}$。注意，$\{GGSW_{\vec{S}_{bk}, \sigma}^{\beta, l}(s_i)\}_{i=0}^{k-1}$ 是使用 GLWE 秘密密钥 $\vec{S}_{bk}$（即\textit{自举}密钥）对 LWE 秘密密钥 $S$ 进行的 GGSW 加密。按照以下流程，使用 $\vec{S}_{bk}$ 同态计算 $V \cdot X^{-\hat{b} + \sum_{i=0}^{k-1}{\hat{a}_is_i}}$，即将 $V$ 的系数向左盲旋转 $\hat{b} + \sum_{i=0}^{k-1}{\hat{a}_is_i}$ 个位置：

\begin{enumerate}
\item 在 TFHE 设置阶段，使用自举密钥 $\vec{S}_{bk}$ 对多项式 $V$ 进行 GLWE 加密，从而使 $V$ 的每个系数都受到加密保护。
\item 计算 $V_0 = V \cdot X^{-\hat{b}}$，本质上就是将 $V$ 的各多项式向左旋转 $\hat{b}$ 个位置。由于 $\hat{b}$ 是 LWE 密文中公开可见的已知值，可以直接计算 $V_0 = V \cdot X^{-\hat{b}}$ 的旋转。
\item 计算 $V_1 = V_0 \cdot X^{\hat{a}_0s_0} = s_0 \cdot (V_0 \cdot X^{\hat{a}_0} - V_0) + V_0$。
该公式适用于 $s_0 \in \{0, 1\}$ 这一特殊情形：若 $s_0 = 0$，则 $V_1 = V_0$；若 $s_0 = 1$，则 $V_1 = V_0 \cdot X^{\hat{a}_0}$。对 $s_0 \cdot (V_0 \cdot X^{\hat{a}_0} - V_0) + V_0$ 的计算通过 TFHE 同态加法和乘法实现。我们将其称为 $V_0$ 的盲旋转，其中选择位 $s_0$（即秘密密钥向量 $S$ 的第一个元素）使用 $\vec{S}_{bk}$（即自举密钥）加密为 GGSW 密文，而 $V_0$ 则是以 GLWE 密文形式加密的多项式。将 GLWE 加密的 $V_0$ 与 $x^{\hat{a}_0}$ 相乘，采用 GLWE 密文与明文乘法（\autoref{sec:glwe-mult-plain}）；从所得结果中减去 GLWE 加密的 $V_0$，采用 GLWE 密文与密文加减法（\autoref{sec:glwe-add-cipher}）；将结果乘以 GGSW 加密的 $s_0$，采用 GLWE 与 GGSW 乘法（\autoref{subsubsec:tfhe-glwe-to-ggsw-multiplication}）；最后加上 GLWE 加密的 $V_0$，采用 GLWE 密文与密文加法。若 $s_0 = 1$，公式中的 $X^{\hat{a}_0}$ 项会乘到 $V_0$ 上，等效于将 $V_0$ 的系数向右旋转 $\hat{a}_0$ 个位置。若 $s_0 = 0$，则 $V_0$ 不发生旋转，保持原样。无论哪种情况，所得 $V_1$ 都以新的 GLWE 密文形式加密。在这一盲旋转过程中，除非知道 $s_0$ 和 $\vec{S}_{bk}$，否则无法判断 $V_0$ 是否发生旋转，也无法知道旋转了多少个位置。
\item 使用与上一步相同的盲旋转方法，计算 $V_2 = V_1 \cdot X^{\hat{a}_1s_1} = s_1 \cdot (V_1 \cdot X^{\hat{a}}_1 - V_1) + V_1$ 的 GLWE 加密。注意，已知的公开信息包括 $\hat{a}_1$、$V_1$ 的 GLWE 加密，以及使用 $\vec{S}_{bk}$ 加密 $s_1$ 得到的 GGSW 密文。
\item 继续以相同方式计算 $V_3, V_4, \gap{$\cdots$}, V_k$ 的 GLWE 加密，最终得到 $V_k$ 的 GLWE 加密；其计算结果等价于：

$V' = V_k$

$\text{ } = V_{k-1} \cdot X^{\hat{a}_{k-1}s_{k-1}}$

$\text{ } = V_0 \cdot X^{\hat{a}_0s_0}X^{\hat{a}_1s_1}\cdots X^{\hat{a}_{k-1}s_{k-1}}$

$\text{ } = V \cdot X^{- \hat{b} + \sum_{i=0}^{k-1}{\hat{a}_is_i}} $

$\text{ } = V \cdot X^{-(\hat{\Delta} m + \hat{e}_b)}$

\end{enumerate}

这意味着，最终多项式 $V_k$ 的 GLWE 加密所对应的常数项系数为 $m$，因为 $V(X)$ 被设计为具有映射（$\hat\Delta m + \hat e_{*}, m$）。

注意，虽然应用所使用的明文空间被限制为 $\mathbb{Z}_t$ 内连续的 $\dfrac{t}{2}$ 个模剩余值（\autoref{subsec:tfhe-zero-padding}），但密文空间没有这一限制。也就是说，盲旋转的中间步骤允许将 $V(X)$ 旋转超过 $n$ 个位置，因为后续步骤可以将无效位置重新转回有效位置。因此，噪声自举的正确性只取决于完成后的状态 $V_k(X)$。旋转完成的 $V_k(X)$ 必须相当于向左旋转了 $\hat\Delta m + \hat e$ 个位置。因此，只须保证 $\hat\Delta m + \hat e$ 落在 $\mathbb{Z}_t$ 域内预先指定的连续 $\dfrac{t}{2}$ 个模剩余值范围中；这等价于确保累计旋转次数至多为 $n-1$ 个位置，以避免提取到符号相互矛盾的系数。

\begin{figure}[h!]
    \centering
  \includegraphics[width=0.2\linewidth]{figures/mux.pdf}
  \caption{MUX 逻辑门示意图}
  \label{fig:mux}
\end{figure}

\para{同态 MUX 逻辑门：} 在步骤 3 中，公式 $s_0 \cdot (V_0 \cdot X^{\hat{a}_0} - V_0) + V_0$ 实现了 \autoref{fig:mux} 所示的 MUX 逻辑门。本例中，$s_0$ 是选择位，用于在两个输入 $V_0$ 和 $V_0 \cdot X^{\hat{a}_0}$ 之间进行选择。如果 $s_0$ 为 1，则输出为 $V_0 \cdot X^{\hat{a}_0}$；否则输出为 $V_0$。在本设计中，同态计算 $s_0 \cdot (V_0 \cdot X^{\hat{a}_0} - V_0) + V_0$ 实际上实现了一个同态 MUX 门，其中两个输入采用 LWE 加密，选择位采用 GGSW 加密。

\para{GLWE 与 GGSW 密文的维数：} 加密 LUT 多项式 $V(X)$ 的 GLWE 密文维数为 $k' \times n$，其中 $k'$ 是 $\vec{A}$ 的长度，即公开掩码多项式的数量，$n$ 是 $A, B,$ 和 $V$ 的多项式次数。加密 $\vec{S}_{\textit{bk}}$ 的每个元素（即自举密钥多项式的每个系数）的 GGSW 密文维数为 $(k' + 1) \times l \times (k' + 1)$。实际中令 $k' = 1$，使这些密文简化为 RLWE 与 RGSW 密文。设置 $k' = 1$ 是出于计算效率的考虑。

\subsubsection{系数提取}

接下来，利用系数提取技术（\autoref{subsec:tfhe-extraction}），提取旋转后的多项式 $V'(X)$' 的常数项系数，得到 $m$ 的加密：$\textsf{LWE}_{\vec{s}', \sigma}(\Delta m)$，其中 $\vec{s}'$ 是长度为 $k'\cdot n$ 的向量。此时，原始密文 $\textsf{LWE}$ 的旧噪声 $\hat{e}_b$ 已经消失，而自举后的新密文 $\textsf{LWE}_{\vec{s}', \sigma}(\Delta m)$ 含有新生成的较小噪声 $e_s$。

事实上，盲旋转过程（\autoref{subsec:bootstrapping-blind-rotation}）中的同态 MUX 逻辑涉及大量密文乘法和加法，直到系数提取之前，这些运算都会累积额外噪声。为限制一系列 MUX 逻辑运算中累积的噪声，可以仔细调整安全参数。例如，可用较窄的高斯分布采样噪声 $e$，同时选取足够大的 $n$，以补偿减小噪声对安全性的影响。另一方面，增大 $q$ 会使系统更容易受到格规约攻击，从而降低安全性。


\subsubsection{密钥切换}
\label{subsec:bootstrapping-key-switch} 

系数提取的输出是 $\textsf{LWE}_{\vec{s}', \sigma}(\Delta m)$，它使用长度为 $k'\cdot n$ 的秘密密钥 $\vec{s}'$ 加密。为保持一致，需要执行 $\vec{s}' \rightarrow \vec{s}$ 的密钥切换，其中 $\vec{s}$ 是原始的长度为 $k$ 的密钥。可采用 \autoref{subsec:tfhe-key-switching} 中小结~\ref*{subsec:tfhe-key-switching} 所介绍的技术完成这一密钥切换。



\subsubsection{噪声自举小结}
\label{subsec:tfhe-summary}

TFHE 的噪声自举过程归纳如下：

\begin{tcolorbox}[title={\textbf{\tboxlabel{\ref*{subsec:tfhe-summary}} TFHE 噪声自举流程}}]

\textbf{\underline{查找表加密}：} 使用长度为 $k'$ 的自举密钥 $\vec{S}_{bk}$，将 LUT 多项式 $V(X)$ 加密为 GLWE 密文。

  
\begin{enumerate}
\item \textbf{\underline{模数切换}：} 对 TFHE 密文 $\textsf{LWE}_{\vec{s}, \sigma}(\Delta m + e_b)$ 执行 $q \rightarrow 2n$ 的模数切换，得到 $\textsf{LWE}_{\vec{s}, \sigma}(\hat{\Delta} m + \hat e_b)$，其中 $\hat\Delta = \Delta \cdot \dfrac{2n}{q} = \dfrac{2n}{t}$。

$ $

\item \textbf{\underline{盲旋转}：} 通过递归计算下式，将 GLWE 加密的多项式 $V$ 向左旋转 $(\hat{b} - \sum\limits_{i=0}^{k-1}{\hat{a}_is_i}) = (\hat{\Delta}m + \hat{e}_b)$ 个位置：

$V_0 = V \cdot X^{-\hat{b}}$

$V_1 = V_0 \cdot X^{\hat{a}_0s_0} = s_0 \cdot (V_0 \cdot X^{\hat{a}_0} - V_0) + V_0$

\text{ } $\vdots$

$V_k = V_{k-1} \cdot X^{\hat{a}_{k-1}s_{k-1}} = s_{k-1} \cdot (V_{k-1} \cdot X^{\hat{a}_{k-1}} - V_{k-1}) + V_{k-1}$

$= V_0 \cdot X^{\hat{a}_0s_0}X^{\hat{a}_1s_1}\cdots X^{\hat{a}_{k-1}s_{k-1}}$

$= V \cdot X^{- \hat{b} + \sum_{i=0}^{k-1}{\hat{a}_is_i}} $

$ = V \cdot X^{-(\hat{\Delta}m + \hat{e}_b)}$

$ $

\noindent 上述盲旋转的每个实际步骤，都通过以下 TFHE 密文与密文乘法和加法计算：

$\textsf{GLWE}_{\vec{S}_{bk}, \sigma}(V_{i+1}) = \textsf{GGSW}_{\vec{S}_{bk}, \sigma}^{\beta, l}(s_i) \cdot (\textsf{GLWE}_{\vec{S}_{bk}, \sigma}(V_i) \cdot X^{\hat{a}_i} - \textsf{GLWE}_{\vec{S}_{bk}, \sigma}(V_i)) + \textsf{GLWE}_{\vec{S}_{bk}, \sigma}(V_i)$

$ $

\item \textbf{\underline{系数提取}：} 从旋转后的多项式 $V_k$ 中同态提取常数项系数 $m$，得到 $\textsf{LWE}_{\vec{s}', \sigma}(\Delta m)$，其中 $\vec{s}'$ 是长度为 $k' \cdot n$ 的向量。

$ $

\item \textbf{\underline{密钥切换}：} 同态执行 LWE 密文的 $\vec{s}' \rightarrow \vec{s}$ 密钥切换。

%\item \textbf{Rescaling:} Multiply $\textsf{LWE}_{\vec{s}, \sigma}(\hat{\Delta}m)$ by $\dfrac{q}{2n}$ to get:  

%$\dfrac{q}{2n} \cdot LWE_{S, \sigma}(\hat{\Delta}m)$

%$\textsf{LWE}_{\vec{s}, \sigma}(\dfrac{q}{2n} \cdot \hat{\Delta}m)$

%$ = LWE_{S, \sigma}(\Delta m)$

%, whose noise get reinitialized to $e_s$. 
\end{enumerate}

$ $

\para{\underline{将可用明文范围减半}：} LUT 多项式 V 的负循环旋转会带来问题。为避免该问题，要求应用确保明文 $m$ 只使用 $\mathbb{Z}_t$ 内连续的 $\dfrac{t}{2}$ 个模剩余值。这样可以防止 $V(X)$ 旋转超过 $n-1$ 个位置，从而避免系数提取得到符号相互矛盾的系数。

$ $

\para{\underline{$k'$ 的选择}：} 为提高计算效率，将 $k'$ 设置为 1，使 GLWE 与 GGSW 密文简化为 RLWE 与 RGSW 密文。
%\item Piping the LUT output to a \textbf{negacyclic function}.
%\item Multiplying the LUT output with \textbf{another negacyclically rotating polynomial}. 
%\end{itemize}
\end{tcolorbox}


\subsubsection{示例：噪声自举}
\label{subsec:tfhe-noise-bootstrapping-ex}

假设 GLWE 的安全参数设置如下：$n = 16$, $t = 8$, $q = 64$, $k = 8$

$\mathbb{Z}_{t=8} = \{-4, -3, -2, -1, 0 , 1, 2, 3\}$

$\mathbb{Z}_{q=64} = \{ -32, -31, -30, \gap{$\cdots$}, 29, 30, 31 \}$

$\Delta = \dfrac{q}{t} = \dfrac{64}{8} = 8$

$ $

并假设有如下 LWE 密文：

$\vec{s} = (1, 0, 0, 1, 1, 1, 0, 1) = \mathbb{Z}_2^{k=8}$

$m = 1  \in \mathbb{Z}_{t=8}$

$\Delta m = 1 \cdot 8 = 8 \in \mathbb{Z}_{q=64}$



$\textsf{LWE}_{\vec{s}, \sigma}(\Delta m) = (a_0, a_1, a_2, a_3, a_4, a_5, a_6, a_7, b) 
= (8, -28, 4, -32, 0, 31, -6, 7, 24) \in \mathbb{Z}_{q=64}^{k+1=9}$

$e = 2 \in \mathbb{Z}_{q=64}$（为保证正确解密，应满足 $|e| < \dfrac{\Delta}{2} = 4$）

$b = \sum\limits_{i=0}^{7}a_is_i + \Delta m + e = (8 -32 + 31 +7) + 8 + 2 = 24 \in \mathbb{Z}_{q=64}$

$ $

此时，TFHE 噪声自举过程如下：

$ $

\begin{enumerate}
\item \textbf{模数切换：} 对 $\textsf{LWE}_{\vec{s}, \sigma}(\Delta m)$ 执行 $q \rightarrow 2n$ 的模数切换，即 $64 \rightarrow 32$。模数切换后，原始 LWE 密文变为：

$\mathbb{Z}_{2n=32} = \{-16, -15, -14, \gap{$\cdots$}, 13, 14, 15\}$

$\vec{s} = (1, 0, 0, 1, 1, 1, 0, 1) = \mathbb{Z}_2^{k=8}$


$\hat{\Delta} = \Delta\dfrac{2n}{64} = 8\dfrac{32}{64} = 4$

$\hat{\Delta}m = 4 \cdot 1 = 4 \in \mathbb{Z}_{2n=32}$

$\hat{e} =  \left\lceil e \dfrac{2n}{q} \right\rfloor = \left\lceil 2\dfrac{32}{64} \right\rfloor = 1 \in \mathbb{Z}_{2n=32}$

$ $

$\textsf{LWE}_{\vec{s}, \sigma}(\hat{\Delta}m) = (\hat{a}_0, \hat{a}_1, \hat{a}_2, \hat{a}_3, \hat{a}_4, \hat{a}_5, \hat{a}_6, \hat{a}_7, \hat{b}) \in \mathbb{Z}_{2n=32}^{k+1=9}$


$ = (\Big\lceil 8\dfrac{32}{64} \Big\rfloor, \Big\lceil -28\dfrac{32}{64}\Big\rfloor, \Big\lceil 4\dfrac{32}{64}\Big\rfloor, \Big\lceil -32\dfrac{32}{64}\Big\rfloor, \Big\lceil 0\dfrac{32}{64}\Big\rfloor, \Big\lceil 31\dfrac{32}{64}\Big\rfloor, \Big\lceil -6\dfrac{32}{64}\Big\rfloor, \Big\lceil 7\dfrac{32}{64}\Big\rfloor, \Big\lceil 24\dfrac{32}{64}\Big\rfloor)$

$ $

$= (4, {-14}, 2, -16, 0, 16, {-3}, 4, 12)$

$ $

注意 $\sum\limits_{i=0}^{7}(\hat{a}_is_i) + \hat{\Delta}m + \hat{e} = (4 - 16 + 16 + 4) + 4 + 1 = 13 \in \mathbb{Z}_{2n=32}$

$ $

$\hat b = 12 \approx 13 = \sum\limits_{i=0}^{7}(\hat{a}_is_i) + \hat{\Delta}m + \hat{e}$

$\hat b$ 中的这一微小差异，来自模数切换时对 $\hat a_0, \hat a_1, \cdots , \hat e$ 取整产生的累积噪声。

$ $

\item \textbf{盲旋转：} 假设应用通过限制可用明文值来避免多项式负循环旋转的问题：仅使用 $\mathbb{Z}_8 = \{-4, -3, -2, -1, 0, 1, 2, 3\}$ 内以下连续的 $\dfrac{8}{2}$ 个模剩余值，即 $\{-2, -1, 0, 1\}$。这意味着 $V(X)\cdot X^{i}$ 中的 $i=\hat\Delta m + \hat e$ 只能取以下值：$i = \{-8, -7, \cdots, 6, 7\}$。根据这些要求，\autoref{tab:lut} 给出了将 $\hat{\Delta} m + \hat{e}$ 映射到 $\Delta m$ 的查找表多项式 $V(X)$。


\begin{table}[h]
\centering
\footnotesize
%\noindent\adjustbox{max width=\columnwidth}{
\begin{tabular}{|c||c|c|c|c|c|c|c|c|c|c|} % left align
\hline
% remove line in a table, \usepackage{adjustbox}
\multicolumn{10}{|l|}{$V(X) = v_0 + v_1X + v_2X^2 + v_3X^3 + v_4X^4 + v_5X^5 + v_6X^6 + v_7X^7$}  \\ 
\multicolumn{10}{|l|}{\textcolor{white}{......} $ + v_8X^8 + v_9X^9 + v_{10}X^{10} + v_{11}X^{11} + v_{12}X^{12} + v_{13}X^{13} + v_{14}X^{14} + v_{15}X^{15}$}  \\ 
\multicolumn{10}{|l|}{\textcolor{white}{....}$= 0 + 0X +0X^2 +0X^3 +1X^4 +1X^5 +1X^6 +1X^7$}  \\ 
\multicolumn{10}{|l|}{\textcolor{white}{......}$+ 2X^8  + 2X^9 + 2X^{10} + 2X^{11} + 1X^{12} + 1X^{13} + 1X^{14} + 1X^{15}$} \\ 
\hline
\hline
\textbf{\boldmath$i = \hat{\Delta}m + \hat{e}$} & $-8$ & $-7$ & $-6$ & $-5$ & $-4$ & $-3$ & $-2$ & $-1$ \\
(in $V \cdot X^{-i}$) & ($\textcolor{orange}{110}\textcolor{green}{00}_2$) & ($\textcolor{orange}{110}\textcolor{green}{01}_2$) & ($\textcolor{orange}{110}\textcolor{green}{10}_2$) & ($\textcolor{orange}{110}\textcolor{green}{11}_2$) & ($\textcolor{orange}{111}\textcolor{green}{00}_2$) & ($\textcolor{orange}{111}\textcolor{green}{01}_2$) & ($\textcolor{orange}{111}\textcolor{green}{10}_2$) & ($\textcolor{orange}{111}\textcolor{green}{11}_2$)\\
\hline
\textbf{常数项的} & $-2$ & $-2$ & $-2$ & $-2$ & $-1$ & $-1$ & $-1$ & $-1$ \\
 \textbf{系数（$V\cdot X^{-i}$）}& $\textcolor{orange}{110}_2$ & $\textcolor{orange}{110}_2$ & $\textcolor{orange}{110}_2$ & $\textcolor{orange}{110}_2$ & $\textcolor{orange}{111}_2$ & $\textcolor{orange}{111}_2$ & $\textcolor{orange}{111}_2$ & $\textcolor{orange}{111}\textcolor{green}{00}_2$ \\
\hline
\textbf{$\bm{m}$（明文）} & $-2$ & $-2$ & $-2$ & $-2$ & $-1$ & $-1$ & $-1$ & $-1$ \\
\hline
\hline
\textbf{\boldmath$i = \hat{\Delta}m + \hat{e}$} & $0$ & $1$ & $2$ & $3$ & $4$ & $5$ & $6$ & $7$ \\
(in $V \cdot X^{-i}$) & ($\textcolor{orange}{000}_2$)& ($\textcolor{orange}{000}_2$)& ($\textcolor{orange}{000}_2$)& ($\textcolor{orange}{000}_2$)& ($\textcolor{orange}{001}_2$)& ($\textcolor{orange}{001}_2$)& ($\textcolor{orange}{001}_2$)&($\textcolor{orange}{001}_2$)\\
\hline
\textbf{常数项的} & $0$ & $0$ & $0$ & $0$ & $1$ & $1$ & $1$ & $1$ \\
\textbf{系数（$V\cdot X^{-i}$）}& $\textcolor{orange}{000}\textcolor{green}{00}_2$ & $\textcolor{orange}{000}\textcolor{green}{00}_2$ & $\textcolor{orange}{000}\textcolor{green}{00}_2$ & $\textcolor{orange}{000}\textcolor{green}{00}_2$ & $\textcolor{orange}{001}\textcolor{green}{00}_2$ & $\textcolor{orange}{001}\textcolor{green}{00}_2$ & $\textcolor{orange}{001}\textcolor{green}{00}_2$ & $\textcolor{orange}{001}\textcolor{green}{00}_2$ \\
\hline
\textbf{$\bm{m}$（明文）} & $0$ & $0$ & $0$ & $0$ & $1$ & $1$ & $1$ & $1$ \\
\hline
\end{tabular}
\centering
\caption{LWE 参数设置为 $n=16, q=64, t=8$ 时的查找表。
\textcolor{orange}{橙色}表示明文 $m$ 的比特。\textcolor{green}{绿色}表示噪声 $e$ 的比特。 %\textcolor{cyan}{Cyan} is the 0-padding bit at the MSB of $\hat\Delta m + e$ (and thus the MSB of $m$), ensured by the application's usage of plaintext numbers.
}
\label{tab:lut}
\end{table}

注意，$V(X)$ 中 $X^8 \sim X^{15}$ 项的系数取 $\{2, 1\}$，而不是 $\{-2, -1\}$。这样，当 $V$ 向左旋转 $\{-8,-7,-6,-5,4,5,6,7\}$ 个位置时，由于越过指数 $n$ 的边界而发生回绕，常数项系数的符号会翻转为 $\{-2, -1\}$。

实际自举时，将 \autoref{tab:lut} 中的 $V(X)$（它已被 GLWE 加密）向左盲旋转 $\hat{b} - \sum\limits_{i=0}^{7}\hat{a}_is_i = 4$ 个位置，计算如下：

$\hat\Delta m + \hat e = \hat{b} - \sum\limits_{i=0}^{7}\hat{a}_is_i = 12 - (4 - 16 + 16 + 4) = 4 \text{ mod 32} \in \mathbb{Z}_{2n=32}$ 


在 \autoref{tab:lut} 中，如果旋转次数 $i = 4$，则对应的常数项系数为 $v_4 = 1 = m$。由于 $\Delta = 4$，最终得到 $\textsf{LWE}_{\vec{s},\sigma}(\Delta m) = 1$。



$ $

实际的盲旋转计算如下：

$\vec{s} = (1, 0, 0, 1, 1, 1, 0, 1)$

$\textsf{LWE}_{\vec{s}, \sigma}(\hat{\Delta}m) = (\hat{a}_0, \hat{a}_1, \hat{a}_2, \hat{a}_3, \hat{a}_4, \hat{a}_5, \hat{a}_6, \hat{a}_7, \hat{b}) =  (4, -14, 2, -16, 0, 16, -3, 4, 12)$

$V_0 = V \cdot X^{-\hat{b}} = V \cdot X^{-12} = v_{12} + v_{13}X + v_{14}X^2 + \cdots $

$V_1 = V_0 \cdot X^{\hat{a}_0s_0} = s_0 \cdot (V_0 \cdot X^{\hat{a}_0} - V_0) + V_0 = V_0 \cdot X^{4} = v_{8} + v_{9}X + v_{10}X^2 + \cdots$ 

$V_2 = V_1 \cdot X^{\hat{a}_1s_1} = s_1 \cdot (V_1 \cdot X^{\hat{a}_1} - V_1) + V_1 = V_1 = v_{8} + v_{9}X + v_{10}X^2 + \cdots$

$V_3 = V_2 \cdot X^{\hat{a}_2s_2} = s_2 \cdot (V_2 \cdot X^{\hat{a}_2} - V_2) + V_2 = V_2 = v_{8} + v_{9}X + v_{10}X^2 + \cdots$

$V_4 = V_3 \cdot X^{\hat{a}_3s_3} = s_3 \cdot (V_3 \cdot X^{\hat{a}_3} - V_3) + V_3 = V_3 \cdot X^{-16} = -v_{8} - v_{9}X - v_{10}X^2 - \cdots$

$V_5 = V_4 \cdot X^{\hat{a}_4s_4} = s_4 \cdot (V_4 \cdot X^{\hat{a}_4} - V_4) + V_4 = V_4 \cdot X^{0} = -v_{8} - v_{9}X - v_{10}X^2 - \cdots$

$V_6 = V_5 \cdot X^{\hat{a}_5s_5} = s_5 \cdot (V_5 \cdot X^{\hat{a}_5} - V_5) + V_5 = V_5 \cdot X^{16} = v_{8} + v_{9}X + v_{10}X^2 + \cdots$

$V_7 = V_6 \cdot X^{\hat{a}_6s_6} = s_6 \cdot (V_6 \cdot X^{\hat{a}_6} - V_6) + V_6 = V_6 = v_{8} + v_{9}X + v_{10}X^2 + \cdots$

$V_8 = V_7 \cdot X^{\hat{a}_7s_7} = s_7 \cdot (V_7 \cdot X^{\hat{a}_7} - V_7) + V_7 = V_7 \cdot X^{4} = v_{4} + v_5X + v_6X^2 + \cdots $

$ $

盲旋转的最终输出是 $V_8$ 的 GLWE 密文 $\textsf{GLWE}_{\vec{S}, \sigma}(V_8)$，其常数项系数为 $v_4 = m = 1$。

$ $

上述盲旋转的每一步都通过以下 TFHE 密文与密文乘法计算：

$\textsf{GLWE}_{\vec{S}, \sigma}(V_{i+1}) = \textsf{GGSW}_{\vec{S}, \sigma}^{\beta, l}(s_i) \cdot (\textsf{GLWE}_{\vec{S}, \sigma}(V_i) \cdot X^{\hat{a}_i} - \textsf{GLWE}_{\vec{S}, \sigma}(V_i)) + \textsf{GLWE}_{\vec{S}, \sigma}(V_i)$

$ $

这一计算留给读者作为练习。

$ $

\item \textbf{系数提取：} 盲旋转结束时，最终得到如下 GLWE 密文：

$\textsf{GLWE}_{\vec{S}, \sigma}(V_8)$

$ = \textsf{GLWE}_{\vec{S}, \sigma}\bm(\hat{\Delta} \cdot (v_4 + v_5X + v_6X^2 + v_7X^3 + v_8X^{4} + v_9X^{5} + v_{10}X^{6} + v_{11}X^{7} + v_{12}X^{8} + v_{13}X^{9} + v_{14}X^{10} + v_{15}X^{11} - v_{0}X^{12} - v_{1}X^{13} - v_{2}X^{14} - v_{3}X^{15})\bm)$

$ = \textsf{GLWE}_{\vec{S}, \sigma}\bm(\hat{\Delta}\cdot(1 + 1X + 1X^2 + 1X^3 + 2X^{4} + 2X^{5} + 2X^{6} + 2X^{7} + 1X^{8} + 1X^{9} + 1X^{10} + 1X^{11} - 0X^{12} - 0X^{13} - 0X^{14} - 0X^{15})\bm)$

$ = \left(A_0 = \sum\limits_{j=0}^{15}(a_{0,0} + a_{0,1}X + \cdots), A_1 = \cdots, A_{k-1} = \cdots, B = \sum\limits_{j=0}^{15}b_{j}X^j\right)$

现在，利用系数提取公式（小结~\ref{subsec:tfhe-extraction}），提取加密多项式 $\textsf{GLWE}_{\vec{S}, \sigma}(\hat{\Delta} \cdot (1 + 1X + 1X^2 + \cdots))$ 的常数项系数。具体而言，提取的是与 $\textsf{LWE}_{\vec{s}, \sigma}(\Delta m_0)$ 对应的常数项系数。通过以下计算提取 $\textsf{LWE}_{\vec{s}, \sigma}(\Delta m_0)$：

$\textsf{LWE}_{\vec{s}, \sigma}(\Delta m_0) = (a_0', a_1', \gap{$\cdots$} , a_{nk-1}', b_h)$ \textcolor{red}{ $\rhd$ 其中 $h = 0$}

\[
    \text{，其中 } a'_{n \cdot i + j} =   
    \begin{cases}
      a_{i,0 - j} \text{（若 } 0 \leq j \leq 0\text{）}\\
      -a_{i,n + 0 - j} \text{（若 } 0+1 \leq j \leq n-1\text{）}\\
  \end{cases}
, b_0 \text{ 取自多项式 } B
\]


\end{enumerate}


\subsubsection{讨论}
\label{subsec:tfhe-noise-bootstrapping-discussion}

\begin{itemize}
\item \textbf{可编程自举}：前述自举（\autoref{subsec:tfhe-noise-bootstrapping}）使用一个简单查找表 $V(X)$，将 $\Delta m + e$ 映射到 $\Delta m$；而通过修改 $V(X)$ 的系数，可以让 $\Delta m + e$ 映射到不同的值。例如，改变 LUT 输入与输出之间的映射，可以实现 AND、OR、XOR、CMUX 等逻辑门，相关内容见 \autoref{subsec:tfhe-noise-bootstrapping-gate}。这种修改 $V(X)$ 中指数与系数之间映射的自举称为可编程自举。如果将 $V(X)$ 加密为 GLWE 密文，就能同时隐藏映射关系及每个输入实例，从而实现\textit{函数加密}。注意，普通自举（\autoref{subsec:tfhe-noise-bootstrapping}）和可编程自举（\autoref{subsec:tfhe-noise-bootstrapping-gate}）产生的噪声量相同。
\item \textbf{自举噪声}：在自举过程中旋转 LUT 多项式 $V(X)$ 时，需要在同态 MUX 门中执行许多 TFHE 乘法以得到 $V_0 \cdots V_k$，因此，在最终完成噪声重新初始化之前，不可避免地会产生额外噪声。不过，仔细选择参数可以限制模数切换和盲旋转期间这些额外噪声的增长。
\end{itemize}


\subsubsection{应用：门自举}
\label{subsec:tfhe-noise-bootstrapping-gate}

除了实现盲旋转中使用的同态 MUX 逻辑门（\autoref{subsec:bootstrapping-blind-rotation}），还可以利用 LUT 多项式 $V(X)$ 实现 AND、NAND、OR、XOR 等其他同态逻辑门。实现这些门时，每个密文加密一个单比特明文；也可以将若干比特通过线性组合公式合并，再用 LUT 同时处理。假设 $q = 32$、$t = 8$、$m \in \mathbb{Z}_8 = \{-4,-3, -2,-1,0,1,2,3\}$、$\Delta = \dfrac{q}{t} = 4$、$\hat\Delta = \dfrac{\Delta\cdot 2n}{q} = 2$，并将逻辑门输入编码为 LWE 明文 $0 \rightarrow -1$ 和 $1 \rightarrow 1$，且最大累积噪声为 $e = [-1, 1]$。


\begin{table}[ht]
  \centering
  \resizebox{\columnwidth}{!}{%
    \begin{tabular}{c c | r r | c | c | c | c}
      \hline
        \multicolumn{8}{c}{查找表多项式 $V(X) = 1 + 1X + 1X^2 + 1X^3 + 1X^4 + 1X^5 + 1X^6 + 1X^7$}\\
      \hline
      \multicolumn{2}{c|}{\textbf{解码值}} & 
      \multicolumn{2}{c|}{\textbf{编码值}} & 
      \multicolumn{1}{c|}{\begin{tabular}{c}
        \textbf{编码的} \\ \textbf{线性} \\
        \textbf{组合}
      \end{tabular}} & 
      \multicolumn{1}{c|}{\begin{tabular}{c}
        \textbf{放大后的} \textbf{编码} \\
         \textbf{组合}
      \end{tabular}} &
      \textbf{自举} & 
      \begin{tabular}{c}
        \textbf{解码} \\ \textbf{结果} \\
      \end{tabular} \\
        \hline
      $d_1$ & $d_2$ & $m_1$ & $m_2$ & $m_1 + m_2 - 1$ & $\hat\Delta\cdot(m_1+m_2-1)$ &
      $V(X)\cdot X^{-\hat\Delta\cdot(m_1 + m_2 - 1)+e}$ &   $d_1 \land d_2$ \\
      \hline
      0 & 0 &
      $-1$ & $-1$ &
      $-3$ &
      $-6$ &
      常数项系数为 $-1$ &
      0 \\
      0 & 1 & 
      $-1$ & $1$ &
      $-1$ &
      $-2$ &
      常数项系数为 $-1$ &
      0 \\
      1 & 0 & 
      $1$ & $-1$ &
      $-1$ &
      $-2$ &
      常数项系数为 $-1$ &
      0 \\
      1 & 1 & 
      $1$ & $1$ &
      $1$ &
      $2$ &
      常数项系数为 $1$ &
      1 \\
      \hline
    \end{tabular}
  }
  \caption{采用额外编码的 AND 运算真值表示例。}
  \label{tab:gate-and}
\end{table}


%\begin{figure}[h!]
%    \centering
%  \includegraphics[width=\linewidth]{figures/torus-gate.png}
%  \caption{An illustration of an AND gate in Torus \href{https://assets-global.website-files.com/622ef9de9152c97467eac748/6298a348a0c75850f0688c83_NAND_bootstrap.png}{(Source)}}
%  \label{fig:torus-gate}
%\end{figure}


\autoref{tab:gate-and} 展示了 AND 逻辑门的可编程自举设计。在这一应用中，将 LUT 多项式 $V$ 定义为 $V(X) = \sum\limits_{i=0}^{7}X^i$。LUT 多项式 $V(X)$ 将一半明文域映射为 $1$，另一半映射为 $-1$（各项越过 $X^7$ 边界时会发生回绕）。在此设计中，每个比特都单独加密为一个独立的 TFHE 密文。门输入 0 和 1 分别编码为 $-1$ 和 $1$。AND 门的线性组合，即同态计算公式，为 $\textsf{LWE}_{\vec{s}, \sigma}(\Delta m_1) + \textsf{LWE}_{\vec{s}, \sigma}(\Delta m_2) - 1$。只有两个输入均为正时，输出才为正（即 $1$）；此时盲旋转将 $V$ 向左旋转 $\hat\Delta\cdot 1 + e$ 个位置，常数项系数为 $1$。因此，盲旋转和系数提取将输出噪声降低后的 $\textsf{LWE}_{\vec{s}, \sigma}(\Delta \cdot 1)$，它是 $1$ 的编码。该设计能够容忍的最大噪声为 $|e| = 1$。若要容忍更大的噪声，应增大 $q$ 和 $n$。

注意，AND 门的 LUT 布局本身具有负循环性质，这是一个特殊情形，因此可以利用 $V(X)$ 固有的负循环旋转性质，将 $V(X)$ 中全部 $2n=16$ 种系数状态用于 AND 门映射函数的输出。然而，在许多应用中，LUT 布局并不一定像该 AND 门示例一样具有负循环性质。前述噪声自举的 LUT 布局（\autoref{subsec:bootstrapping-overview}）也不是负循环函数，而是恒等函数（因为它仅仅移除噪声）。因此，对于大多数应用，需要只使用 $t$ 明文空间中的 $\dfrac{t}{2}$ 部分，以避免 $V(X)$ 的旋转超过 $n-1$ 个位置（\autoref{subsec:tfhe-zero-padding}）。

除 AND 门之外，其他逻辑门也可以按类似方式构造，每种门采用不同的线性组合公式和 LUT 布局。

\para{除法：} TFHE 不支持任意大小明文数的直接除法。这是因为 TFHE 的 LWE 向量元素属于环 $Z_q$，其中每个元素 $g$ 不一定存在乘法逆元 $g^{-1}$，因而难以通过将目标数乘以 $g^{-1}$ 来实现除法。可以改用门自举所实现的逻辑门来完成二进制除法。为支持二进制除法，每个明文必须是一个比特，并被加密为独立的密文。也可以设计线性组合公式，将多个比特合并并同时处理，类似于之前为处理 AND 门的两个输入比特所设计的线性组合。

\subsubsection{应用：神经网络中的自举}
\label{subsec:tfhe-neural-network}

\begin{figure}[h!]
    \centering
  \includegraphics[width=0.8\linewidth]{figures/neural-network.pdf}
  \caption{神经网络示意图}
  \label{fig:neural-network}
\end{figure}

\begin{figure}[h!]
    \centering
  \includegraphics[width=0.8\linewidth]{figures/nn-homomorphic.pdf}
  \caption{神经网络中的可编程自举示意图}
  \label{fig:neural-network2}
\end{figure}

同态加密可以应用于深度神经网络中的神经元。通常，每个神经元的计算包含两个步骤：

\begin{enumerate}
  \item \textbf{输入值的线性组合：} 给定输入特征值（或中间值）集合 $(x_1, x_2, \gap{$\cdots$}, x_n)$、权重集合 $(w_1, w_2, \gap{$\cdots$}, w_n)$ 和偏置 $b$，计算：$y = \sum\limits_{i=1}^{n}x_iw_i + b$。
\item \textbf{激活函数：} 计算 $f(y)$，其中 $f$ 是非线性激活函数，例如 $\sin$ 函数、ReLU、sigmoid、双曲正切等。
\end{enumerate}

TFHE 可以将第一步的线性组合公式 $y = \sum\limits_{i=1}^{n}x_iw_i + b$ 同态计算为 $\sum\limits_{i=1}^{n}\textsf{LWE}_{\vec{s}, \sigma}(x_i) \cdot w_i + b$，其实现使用密文加法（\autoref{sec:glwe-add-cipher}）以及密文与明文乘法（\autoref{sec:glwe-mult-plain}）。

然而，第二步的非线性函数无法直接表示为密文的加法和乘法。为解决这一问题，可以通过可编程自举求值激活函数，使自举的输出与激活函数的输出相同或接近。如果在第二步使用自举，则可以在每个神经元计算结束时刷新噪声，因此原则上能够处理任意深度的神经网络，而无须担心噪声持续增长。



\subsection{离散环面上的 TFHE}
\label{subsec:torus}

\textbf{- 参考资料：}
\href{https://eprint.iacr.org/2021/1402.pdf}{（离散化）环面上的
全同态加密
指南}~\cite{torus}

$ $

环面 $\mathbb{T}$ 是介于 0 和 1 之间并按模 1 回绕的连续实数域，即 $[0, 1)$。

离散环面 $\mathbb{T}_t$ 是如下有限实数集合：$\Big(0, \dfrac{1}{t}, \dfrac{2}{t}, \gap{$\cdots$}, \dfrac{t - 1}{t}\Big)$


前面各小节介绍的 TFHE 方案采用如下设置：

$ $

$m \in \mathbb{Z}_t$ 
$\vec{s} = \{0, 1\}^{k}$
$e \in \mathbb{Z}_q$

$\textsf{LWE}_{\vec{s}, \sigma}(\Delta m) = (a_0, a_1, \gap{$\cdots$}, a_k, b) \in \mathbb{Z}_t^{k + 1}$

$ $

而原始 TFHE 方案是基于离散环面设计的：


$ $

$m \in \mathbb{T}_{t}, \text{ } \vec{s} \in \{0, 1\}^{k}, \text{ } e \in \mathbb{T}_{q}$

$ $

$\textsf{LWE}_{\vec{s}, \sigma}(m) = \textsf{ct} = (a_0, a_1, \gap{$\cdots$}, a_k, b) \in \mathbb{T}_{q}^{k+1}$

$b = \sum\limits_{i=0}^{k}(a_is_i) + m + e \in \mathbb{T}_q$

$\textsf{LWE}_{\vec{s}, \sigma}^{-1}(\textsf{ct}) =  \Big\lceil b - \sum\limits_{i=0}^{k-1}(a_is_i) \Big\rfloor_{\frac{1}{t}} = \Big\lceil m + e \Big\rfloor_{\frac{1}{t}} = m$，其中给定 $e < \dfrac{1}{2t}$

\textcolor{red}{ $\rhd$ 其中 $\lceil x \rfloor_{\frac{1}{t}}$ 表示将 $x$ 舍入为最接近的 $\dfrac{1}{t}$ 的整数倍}

$ $

原始 TFHE 的区别在于，所有值（多项式系数或向量元素）都在模 1 的浮点数域（即 $[0, 1)$）中计算，而不是在大整数域（即 $[0, q)$）中计算。这意味着明文也必须编码为 $[0, 1)$ 中的值，而不是 $[0, q)$ 中的整数。注意，原始 TFHE 方案不需要尺度因子 $\Delta$，因为环面 $[0, 1)$ 的连续域提供了可区分 $q$ 个离散小数值的浮点精度；只要 $E$ 中的每个系数（或向量元素）$e_i$ 小于 $\dfrac{1}{2t}$，解密过程就可以成功消除噪声 $E$。

基于环面和基于整数环的 TFHE 方案均建立在相同的基本原理之上。
